\documentclass[11pt,letterpaper]{article}

\usepackage{fontspec}
\usepackage{amsmath}
\usepackage[margin=0.85in,headheight=22pt,includeheadfoot]{geometry}
\usepackage{microtype}
\usepackage{graphicx}
\usepackage{pdflscape}
\usepackage{caption}
\usepackage{xcolor}
\usepackage{booktabs,tabularx,longtable,array}
\usepackage{enumitem}
\usepackage{fancyhdr}
\usepackage{lastpage}
\usepackage{setspace}
\usepackage{xurl}
\usepackage[hidelinks,pdfauthor={Linkwise Project},pdftitle={Linkwise: An Architecture for AI-Assisted HF Link Management}]{hyperref}

\definecolor{ovnavy}{HTML}{16324F}
\definecolor{ovblue}{HTML}{2E6F95}
\definecolor{ovteal}{HTML}{2A9D8F}
\definecolor{ovgold}{HTML}{E9C46A}
\definecolor{ovpale}{HTML}{EFF5F8}
\definecolor{ovgray}{HTML}{52606D}
\definecolor{ovred}{HTML}{B23A48}

\setlength{\parindent}{0pt}
\setlength{\parskip}{5pt}
% Allow modest extra interword space on difficult lines before overflowing.
\setlength{\emergencystretch}{1em}
\setlist[itemize]{leftmargin=1.35em,itemsep=2pt,topsep=3pt}
\setlist[enumerate]{leftmargin=1.6em,itemsep=3pt,topsep=4pt}
\setlength{\headwidth}{\textwidth}
\pagestyle{fancy}
\fancyhf{}
\lhead{\small\color{ovgray}LINKWISE \textbullet{} SYSTEM ARCHITECTURE}
\rhead{\small\color{ovgray}Whitepaper 0.2 \textbullet{} Proposed}
\cfoot{\small\color{ovgray}\thepage\ of \pageref{LastPage}}
\renewcommand{\headrulewidth}{0.4pt}
\renewcommand{\headrule}{\hbox to\headwidth{\color{ovblue}\leaders\hrule height \headrulewidth\hfill}}
\setcounter{secnumdepth}{2}
\setcounter{tocdepth}{2}

\newcommand{\system}{\textsc{Linkwise}}
\newcommand{\todo}{\textcolor{ovred}{\textbf{TBD}}}
\newcommand{\evidence}[1]{\textsf{#1}}
\newcommand{\urlref}[2]{\href{#1}{#2}}
% Every PlantUML figure gets a dedicated landscape page, a bounded image box,
% and a caption/label. Keep height below the landscape text area for captions.
\newcommand{\landscapeplantuml}[3]{%
  \clearpage
  \begin{landscape}
    \thispagestyle{plain}
    \centering
    \includegraphics[width=0.94\linewidth,height=0.86\textheight,keepaspectratio]{#1}\par
    \vspace{0.10in}
    \captionsetup{width=0.94\linewidth}
    \captionof{figure}{#2}
    \label{#3}
  \end{landscape}
  \clearpage
}

\begin{document}

\hypersetup{pageanchor=false}
\begin{titlepage}
  \pagecolor{ovpale}\color{ovnavy}
  \vspace*{0.6in}
  {\Large\bfseries SYSTEM ARCHITECTURE WHITEPAPER}\par
  \vspace{0.35in}
  {\Huge\bfseries Linkwise}\par
  \vspace{0.16in}
  {\LARGE AI-Assisted HF Link Management with VARA and Winlink}\par
  \vspace{0.35in}
  {\large Software architecture, operating concept, use cases, and design basis}\par
  \vspace{0.65in}
  \vfill
  \textbf{Document:} LINKWISE-WP-001\par
  \textbf{Version:} 0.2 --- Architecture proposal\par
  \textbf{Date:} November 2026\par
  \textbf{Status:} Proposed system; radio MCP and native modem prototypes; operational qualification pending\par
  \vspace{0.25in}
  {\small\color{ovgray}Prepared as an engineering whitepaper. All protocol, automation, and propagation capabilities remain subject to implementation and validation.}
\end{titlepage}
\hypersetup{pageanchor=true}
\nopagecolor
\color{black}

\landscapeplantuml{build/diagrams/Linkwise_Cover_Architecture}{C4 system context: Linkwise and its people, station interfaces, Winlink services, and evidence sources.}{fig:c4-system-context-cover}

\begin{abstract}
\noindent\system{} is a proposed system for helping a radio station send and receive Winlink messages over HF radio. It brings together radio settings, the VARA modem, Winlink gateways, and information about radio conditions. An AI agent would use this information to choose a promising gateway, configure the radio, and manage message sessions automatically within limits set before operation. The operator would not need to supervise routine activity and could stop transmissions using an independent emergency stop. The system would record session results to guide future choices. VOACAP forecasts can help compare possible radio paths, but they cannot guarantee a connection; the possible effects of satellite reentry remain a research question. One radio can use only one HF channel at a time. Linkwise has not yet implemented this full system, so the paper describes a design that still needs to be built and tested.
\end{abstract}

\vspace{0.12in}
\begin{center}
\colorbox{ovpale}{\parbox{0.94\textwidth}{\textbf{Design principle.} Let the agent manage radio configuration, PTT, and link sessions autonomously within a preinstalled station authority. Keep an independent human emergency stop; expose no other runtime operator actuator.}}
\end{center}

\tableofcontents
\newpage

\section{Executive summary}

Reliable HF messaging depends on several things lining up: a gateway must be available, the path must support the selected frequency, the station must be configured correctly, and the modem and client must complete a session. Those facts are spread across radio-control software, VARA, Winlink, channel directories, propagation services, and recorded station experience.

\system{} proposes a local autonomous coordination layer that brings those inputs together. It maintains a time-stamped set of eligible gateway channels, ranks them for the installed service objective, configures the radio, and starts an eligible session. Before each transmission, deterministic policy checks station limits and issues a bounded authorization automatically. Afterward, the system records outcomes to improve future choices.

The design keeps each responsibility at a clear boundary. Hamlib/rigctld exposes all capabilities reported by the radio backend, and the AI can request any locally enabled CAT setting or registered device extension through typed MCP tools. PTT state, key-up, and key-down are also explicit tools; key-up remains subject to station-configured authority, and release remains an independent safety operation. A modem adapter handles protocol timing and the data link, and a Winlink-client adapter manages message and gateway sessions. Other adapters collect propagation and gateway evidence. The model does not connect to CAT/PTT sockets itself: Linkwise serializes and audits each operation and validates it against the radio profile and policy.

Linkwise remains a proposed architecture. Companion architecture, requirements,
interface, safety, simulation, and lab-notes documents carry detailed design
and implementation history.

\section{Purpose and scope}

\subsection{Purpose}
This whitepaper explains the proposed architecture, how the autonomous agent uses the existing radio and messaging software, how the station owner commissions it before deployment, and how the design could be verified. It is written for radio operators, software architects, modem developers, radio integrators, and AI-host developers. It includes a summarized layered SDD view; the companion SDD remains the contract-level design baseline, and this paper's coverage and remaining gaps are assessed in Appendix~\ref{app:template-alignment}.

\subsection{System goals}
Linkwise is intended to:
\begin{itemize}
  \item maintain an accurate, time-bounded set of eligible Winlink RMS channels;
  \item autonomously select channels that balance predicted reachability, live observations, station fit, service status, and measured history;
  \item coordinate radio, modem, and client state without conflicting controllers;
  \item autonomously configure the radio, authorize and conduct bounded sessions, and recover/retry without routine operator intervention;
  \item explain why a gateway was selected and what evidence supports the decision;
  \item improve future selection from actual success, failure, time, and throughput data; and
  \item provide a human-only emergency stop independent of AI, application, and network availability, and fail safely on faults.
\end{itemize}

\subsection{Non-goals}
Linkwise will not replace the Winlink CMS or gateway network, promise worldwide coverage, or claim that a propagation model can predict an exact live opening. It will not turn unverified observations into modem protocol facts. One radio cannot hold simultaneous HF links, and coordinating multiple radios is outside the initial scope.

The following sections establish the communications context, define the link and session terms, and explain the evidence and scoring basis before describing the operating concept.

\section{Background: the communications software}
Readers do not need prior experience with amateur-radio software to follow
the design. A station that sends a Winlink message over HF typically combines
several programs and hardware interfaces, each responsible for a different
part of the work. Linkwise is proposed as a coordinator around those parts,
not a replacement for the radio, the Winlink network, or the existing VARA
modem.

\subsection{The radio and CAT control}
The transceiver is the RF endpoint: it selects a frequency and mode, receives
signals, and transmits audio as radio-frequency energy. Many radios expose a
serial or USB command interface commonly called CAT (computer-aided
transceiver control). CAT is not one universal wire protocol; command details
and available controls vary by manufacturer and model. A software controller
must therefore know which operations a particular radio supports and confirm
state by reading back the radio where possible.

Hamlib is an established open-source library that gives applications a
common programming interface to many radio models. Its command-line client,
\texttt{rigctl}, is useful for interactive control and testing. Its
\texttt{rigctld} daemon exposes radio-control operations over a TCP service,
so separate applications can communicate with Hamlib without embedding its
C library directly [1,21]. Linkwise proposes to use a local Hamlib/\texttt{rigctld}
adapter as the radio-control boundary: Linkwise discovers the active backend's
capabilities, converts them into typed internal operations, serializes writes,
and verifies readback. CAT operations can include more than frequency and
mode—where the selected backend reports support, controls may include VFO,
filter, levels, functions, split, and other radio settings. The available set
is radio-specific, not a promise that every radio implements every operation.

PTT (push-to-talk) is the act of keying the transmitter. Depending on the
station, PTT may be controlled through CAT, a serial control line, a USB audio
interface, or a modem/client integration. It is related to, but distinct from,
the radio's frequency and mode settings. Linkwise therefore models PTT as a
separate safety-critical path: the modem can request transmit timing, while a
local arbiter enforces session authorization, duration limits, release, and
the independent human emergency stop. A radio-specific PTT route must be
qualified; the architecture does not assume all radios behave alike.

\subsection{The modem and audio path}
A sound-card data modem turns digital data into audio waveforms for the
transmitter and decodes received audio back into data. The station connects
the computer's audio input/output to the transceiver, often through a USB
audio interface or an integrated radio codec. A separate PTT route keys the
radio at the right time. This division explains why a working CAT connection
alone is not enough: frequency/mode control, audio routing, and PTT all have
to agree for a session to work.

VARA HF is an existing third-party HF modem used with Winlink and other
applications. The Winlink project distributes VARA products but describes
them as third-party software, not software managed by the Winlink team [22].
In the proposed deployment, VARA remains an external modem process. Winlink
setup material identifies local TCP command/data interfaces for VARA HF
(commonly ports 8300/8301, subject to local configuration) [29]. A Linkwise
adapter could use that client-facing boundary to coordinate modem lifecycle
and session state, route audio, and request PTT through the arbiter. The exact
commands, status fields, licensing, and automation behavior must be verified
against the selected VARA release before they are treated as stable
contracts. Linkwise does not claim to implement or replace
the VARA waveform. A future native modem is a separate research and
interoperability effort.

The initial target transceiver is the Icom IC-7200. Its USB connection
provides CI-V computer control and audio transfer, so the radio adapter can
use Hamlib/rigctld for supported CAT state and controls while the audio
adapter routes receive and transmit audio to the selected modem [45, 46].
USB receive audio is the radio's demodulated audio output; it is not a
wideband RF/IQ stream and therefore cannot survey other channels while the
IC-7200 is tuned to an active session. Such band-level observation requires
a separate receive-only SDR. The IC-7200's CAT meter readings, where the
installed Hamlib backend exposes them, are useful local measurements but
must not be treated as proof of remote protocol reception or link success.

\subsection{Winlink clients, RMS gateways, and message services}
Winlink is a radio-email and messaging system, not a single radio protocol.
At a client station, a Winlink application manages the outbox; messages may
be composed by a user or supplied by an upstream application. The client
selects or is configured for a connection method. For an HF radio session,
the client works with a modem such as VARA HF and a remote RMS (Radio Message
Server), which provides the radio gateway. RMS stations connect messages
into the wider Winlink system through CMS (Common Message Servers); the
network also has
store-and-forward capabilities for some operating conditions [11]. These
roles matter because a successful RF connection is not identical to a
completed message transfer, and a transferred message is not necessarily
proof of final recipient delivery.

Winlink publishes live RMS channel information, including gateway identity,
frequencies, service/status information, and report age [3]. Linkwise can
use that directory to maintain candidate gateway channels, but the directory
is not a live propagation sensor and does not establish that a particular
station can be reached from the local radio. The proposed Winlink client
adapter coordinates a session and collects only fields available through a
supported client interface. Queue counts and transfer progress are optional;
if a client exposes no documented interface for them, Linkwise reports that
capability as unavailable instead of scraping a graphical interface.

\subsection{Propagation prediction and current evidence}
VOACAP (Voice of America Coverage Analysis Program) estimates HF circuit
performance from path endpoints, frequency, time, and station assumptions.
It helps answer a planning question such as “which frequencies and hours are
statistically promising for this path?” It does not determine whether an RMS
gateway is currently online, whether local interference is present, or
whether the next VARA session will succeed. Linkwise therefore uses VOACAP as
one time-bounded planning input, alongside RMS status, observations, and its
own measured session outcomes—not as a transmit authorization or a guarantee
of reachability [2].

ALE systems provide a useful evidence pattern for measured link quality:
stations transmit known soundings, receivers associate decoded measurements
with the sounding station and channel, and systems may average, age, exchange,
and threshold those observations before automatic channel selection [47,
48]. Linkwise may borrow these data-management concepts, but the initial
system does not implement MIL-STD-188-141 ALE and must not label unrelated
VARA or IC-7200 measurements as ALE LQA. Its RF evidence is instead separated
by source and meaning: SDR-derived spectrum/occupancy observations describe
local channel conditions; radio meter samples describe the IC-7200 receiver;
modem events establish protocol-level peer responses; and Winlink client
events establish session and transfer progress.

For wider RF-environment sensing, a configured receive-only SDR can feed a
GNU Radio flowgraph using SoapySDR-supported hardware [49, 50]. The flowgraph
may produce bounded measurements such as noise floor, occupied bandwidth,
energy-detection candidates, and observation coverage. These are local
observations, not decoded station identities or a reachability guarantee.
Raw IQ remains inside the SDR/DSP pipeline; Linkwise ingests compact,
timestamped results through a local adapter. The adapter records receiver,
frequency span, sample/configuration metadata, units, calibration status,
source and collection times, and quality flags. Continuous IQ and waterfall
frames are not MCP payloads or ordinary audit-log entries.

\subsection{MCP and the role of the AI agent}
MCP (Model Context Protocol) is a standard interface through which an AI host
can discover and call typed tools and read contextual resources from a
software server [5]. It is the agent-facing interface, not the radio-control
protocol and not the AI model itself. Linkwise proposes an MCP server that
exposes tools such as reading station state, listing eligible RMS channels,
requesting a radio configuration, starting a bounded session, and observing
its result. The AI host can use those tools to plan and coordinate actions;
Linkwise's deterministic application and policy services still validate and
execute every side effect. The model cannot connect directly to CAT, PTT,
VARA, or Winlink sockets or expand its preinstalled authority.

\subsection{How the existing pieces become a communications solution}
Taken together, the interfaces form a chain rather than one monolithic
program. A local Linkwise service reads the radio through Hamlib, obtains
gateway candidates from Winlink's published channel information, and gathers
propagation and station evidence. Its ranker chooses an eligible next RMS
channel for the installed service objective. The station actor tunes the
radio and verifies readback. The Winlink client coordinates the message
session with VARA HF; VARA handles the modem waveform over the audio path and
requests PTT through Linkwise's safety arbiter. The client and modem report
session events, which Linkwise records and uses to update measured gateway
history. The AI participates through MCP throughout this loop, while local
policy retains the authority to deny unsafe or out-of-profile operations.

This decomposition is also the integration strategy. Existing programs stay
behind adapters, so their version-specific interfaces can be tested and
replaced without putting those details into the optimizer. The current Rust
prototype uses the official Rust MCP SDK \texttt{rmcp} and connects to
Hamlib's \texttt{rigctld} service over TCP. It does not yet include audio
support or embed Hamlib's C library. The cross-platform audio I/O library
\texttt{cpal} remains an option for later evaluation, not a selected
dependency [23,24]. The architecture favors connecting to existing local
services over embedding native libraries until packaging, platform coverage,
licensing, error handling, and hardware behavior have been evaluated.

\subsection{Stakeholders and concerns}
The design serves people with different priorities. The station owner needs reliable pre-deployment configuration and a dependable emergency stop; developers need stable interfaces and testable behavior; safety reviewers need to know how autonomous actions are bounded. The table connects those concerns to the views used later in the paper.

\clearpage
\begin{landscape}
\thispagestyle{plain}
\captionof{table}{Stakeholders, primary concerns, and the architecture views that address them.}
\label{tab:stakeholders}
\begin{tabularx}{\linewidth}{@{}p{1.8in}X p{2.0in}@{}}
\toprule
\textbf{Stakeholder} & \textbf{Primary concern} & \textbf{View covering it}\\
\midrule
Radio operator & Timely message delivery, understandable choices, and independent emergency stop & Operational; safety \\
Station maintainer & Reliable setup, diagnostics, recovery, radio compatibility & Deployment; interfaces \\
Modem developer & Evidence-based interoperability and reproducible tests & Functional; verification \\
Winlink/client integrator & Correct gateway/channel handling and supported client interfaces & Functional; data \\
AI-host integrator & Discoverable tools, clear errors, bounded permissions & Functional; security \\
Safety/security reviewer & Prevention of unintended transmission and auditable decisions & Safety/security \\
Researcher & Evidence provenance and a fair test of propagation hypotheses & Data; simulation \\
\bottomrule
\end{tabularx}
\end{landscape}
\clearpage

\section{Link and session model}
Linkwise tracks a route from discovery through application outcome. The stages
are deliberately separate: an RMS listing is inventory evidence; a candidate
is an eligible station/channel pairing; a ranking is a time-bounded comparison;
preparation configures and verifies the station; authorization is a local
policy decision; an attempt is one bounded session; a modem link is a protocol
connection; and a Winlink transfer is a client-level result. CMS acceptance
and recipient delivery are recorded only when the client explicitly confirms
them. These distinctions define what each observation can establish.

\begin{table}[htbp]
\centering
\caption{Core Linkwise terms and the evidence each term represents.}
\label{tab:linkwise-terms}
\begin{tabularx}{\linewidth}{@{}p{1.35in}X@{}}
\toprule
\textbf{Term} & \textbf{Meaning} \\
\midrule
Gateway channel & One published RMS station, frequency, protocol, and supported mode record. \\
Candidate & A gateway channel that passes the installed station and service eligibility checks. \\
Ranking & A versioned comparison of eligible candidates for an objective and evidence snapshot. \\
Prepared route & A candidate whose radio settings and required subsystem readiness have been checked. \\
Attempt & One authorized, bounded execution of a Winlink session flow. A retry is a new attempt. \\
Modem link & A protocol-level connection reported by the modem. It does not prove client transfer. \\
Transfer outcome & A client-confirmed session or message-transfer result; CMS acceptance and recipient delivery remain distinct when reported. \\
\bottomrule
\end{tabularx}
\end{table}

\section{Propagation and evidence foundations}
\subsection{Propagation physics and measurement limits}
The literature supports a causal, layered account of HF link conditions:
solar and geomagnetic forcing changes ionospheric state; the state controls
wave refraction, absorption, and scattering; path effects combine with
antennas, noise, and equipment to determine the received channel; the modem
and Winlink client then determine whether that channel supports a session and
transfer. These layers have different observables and uncertainty. The
architecture therefore stores propagation estimates separately from measured
radio contact, modem connection, client acceptance, and delivery.

At first order, the electron plasma frequency is
\[
  f_p = \frac{1}{2\pi}\sqrt{\frac{N_e e^2}{\epsilon_0 m_e}},
\]
where \(N_e\) is electron density, \(e\) is electron charge, \(m_e\) is
electron mass, and \(\epsilon_0\) is free-space permittivity. In an
unmagnetized, collisionless plasma the refractive index obeys
\(n^2 \simeq 1-f_p^2/f^2\). This approximation explains why electron density
and operating frequency govern refraction and a vertical-incidence critical
frequency. The terrestrial ionosphere is magnetized, collisional, spatially
structured, and time varying; practical HF analysis therefore requires
magnetoionic dispersion (with ordinary and extraordinary modes), path geometry,
and absorption treatment rather than applying the simple expression as a
complete prediction [31,32,34]. Ray tracing is a computationally useful
geometric-optics approximation where the medium varies slowly relative to the
wave; full-wave methods are needed where that approximation or assumed layer
geometry is inadequate [31].

The complex refractive index also encodes attenuation. Absorption depends on
electron density and collision frequency along the traversed path, with
important contributions from the lower ionosphere. Zawdie et al. show that
absorption calculations depend on the dispersion formulation and the
collision-frequency model; under their comparisons, selecting consistent
collision-frequency assumptions brought the Appleton--Hartree and Sen--Wyller
calculations to within a few percent in the D/E regions, while mismatched
assumptions produced much larger differences [40]. This result is a warning
against treating model names as interchangeable or attaching confidence to a
forecast without recording its ionospheric and absorption assumptions.

Circuit quantities must also be distinguished. A critical frequency
characterizes vertical-incidence reflection under specified conditions; the
maximum usable frequency is path- and mode-dependent. The lowest usable and
optimum working frequencies also depend on losses, noise, and the desired
reliability. A field-strength or SNR estimate is not itself a probability of
successful VARA/Winlink delivery. The latter also depends on time-varying
fading, multipath delay and Doppler, interference, receiver performance,
modem mode, and client behavior [32,34--36]. The link utility model must
therefore keep modeled circuit support and measured service outcomes as
separate response variables.

Measurement types observe different parts of this chain. A vertical ionogram
probes the ionosphere near the sounder; an oblique sounder or calibrated
beacon is more directly associated with an actual radio path; wideband
soundings can estimate delay/Doppler spreading and the channel scattering
function [35,36]. GNSS TEC is an electron-density integral along a GNSS ray,
not an HF-frequency electron-density profile or a measurement of a particular
HF circuit [33,34]. Amateur-radio reporting networks provide valuable,
time-stamped observations of real operating paths and have been used in
peer-reviewed studies of flares and traveling ionospheric disturbances
[41,43]. Their reports are opportunistic and mode/population dependent;
researchers describe them as networks designed to reveal propagation
openings, rather than finely calibrated ionospheric instruments [42]. Thus,
an observed report is evidence that a particular transmission was received;
no report is ambiguous and must not be recorded as proof of path closure.

The resulting evidence-to-decision contract is:
\begin{enumerate}
  \item Store source value, physical unit, source/model version, issue and
        observation times, validity window, path/location, frequency, and
        assumptions. Keep forecast output, direct measurement, reception
        report, and session outcome as different evidence classes.
  \item Convert each input to a stated path/time/frequency feature, with
        missingness, freshness, spatial coverage, uncertainty, and quality
        preserved. Do not treat missing records or no spots as zero signal or
        as a failed path observation.
  \item Compare forecasts against measurements of the same target and scale.
        Evaluate field strength, channel availability, and application-level
        transfer success separately; do not validate an SNR prediction only
        against session success or call both results “propagation accuracy.”
  \item Use model estimates as calibrated statistical evidence for ranking,
        not as a live reachability assertion or authorization input. Retain
        exact input/model IDs so each recommendation and later outcome can be
        replayed.
\end{enumerate}

This synthesis establishes the theory and engineering basis for the
capability, but it does not qualify an implementation. Remaining empirical
work is to characterize the chosen local VOACAPL/VOACAP configuration, build
a path-matched dataset with calibrated observations and station-specific
outcomes, test alternative model and assimilation approaches on
time/path/season-held-out data, and quantify the improvement in the link
ranking objective. The evaluation should report bias and error for continuous
estimates, calibration and discrimination for probabilities, and
denominator-bearing outcome and utility measures for ranking [44]. Any experimental
covariate, including reentry-related evidence, remains shadow-only until it
shows repeatable out-of-sample value over the non-experimental baseline.


\section{Gateway scoring and continuous management}
\subsection{Evidence-to-score model}
For candidate gateway \(g\), the ranker estimates utility as
\[
U(g,t)=P_g(\text{delivery within budget}\mid E_t)\,V_g
       -C_g^{\mathrm{time}}-C_g^{\mathrm{switch}}-C_g^{\mathrm{risk}},
\]
where \(E_t\) is timestamped evidence and \(V_g\) represents user-defined delivery value. Inputs can include RMS availability, path-specific VOACAP results, current reported activity, local receiver observations, antenna bearing, modem compatibility, and historical session outcomes. Any learned estimate must disclose calibration limits and uncertainty.

\subsection{Evidence sources and operational interpretation}
The planner treats a ``channel'' as an eligible RMS gateway/channel record (including its published frequency and supported mode), not a newly created Winlink service or a spectrum allocation. For each candidate path, the evidence service supplies the station and gateway locations, UTC planning interval, frequency, antenna/power assumptions, and compatible mode to a pinned local VOACAPL/VOACAP implementation. The result is stored with those inputs and the model version. The ranker uses the relevant hourly/monthly circuit reliability and signal-to-noise outputs to establish a statistical prior and identify promising frequency/time windows; it can combine that prior with current RMS status, reported activity, local measurements, and empirical outcomes. It cannot claim that the path is open now. Before an attempt, the link manager refreshes candidate inventory and station state, applies the installed station profile and regulatory limits, then autonomously selects and tunes the best eligible listed route for the bounded session flow.

Reentry evidence follows a distinct path. A provider such as KeepTrack's v4 endpoint reports upcoming reentry predictions derived from Space-Track TIP messages, with updated predictions deduplicated by object [18]. These are forecasts with uncertainty, not confirmed event times or direct measurements of ionization. The research pipeline records object identifier, prediction issue time, predicted time/window, source, event status, and available spatial metadata. It aligns the event interval with geographically and temporally relevant sporadic-E/ionosonde observations and Linkwise's measured per-path session results. A reentry-associated feature is therefore a candidate explanatory variable for local E-region variability—not a universal channel-quality adjustment.

The production ranker initially excludes this feature. In shadow/research mode, the system computes a parallel experimental score and explanation, then evaluates it retrospectively and prospectively against (a) a VOACAP-plus-operational-evidence baseline and (b) the same baseline without reentry features. Evaluation uses event/control windows, geographic relevance, source/prediction age, sufficient repeated observations, calibration, and held-out time/path scenarios. Only after controlled testing demonstrates repeatable out-of-sample improvement may a newly reviewed, calibrated profile be installed before a later deployment; enabling it is not a runtime operator action. Its weight and provenance remain visible, and it can never override eligibility, channel-clear checks, local policy, or PTT authorization. This conservative gate is warranted because recent peer-reviewed work examines a possible association between artificial reentries, metallic ions, and E-region anomalies, but does not establish a real-time HF path predictor [19].

VOACAP reliability is statistical for modeled hour/month and is not a guarantee for a particular session [2]. RMS status is a reported gateway state, not RF reachability [3]. PSK Reporter and DX Maps provide reported activity and can be biased by who is transmitting/reporting. NOAA SWPC products such as solar radio flux, geomagnetic indices and forecasts, alerts, and D-Region Absorption Prediction (D-RAP) can characterize space-weather conditions that may affect HF circuits; they are contextual or mechanism-specific evidence, not end-to-end path measurements [24,27]. NASA CCMC DONKI contributes event analyses and linked solar/heliospheric context, while NASA GNSS-derived TEC products may support regional ionospheric studies. DONKI is research/prototyping-quality information, and TEC is not a direct measure of HF circuit success [25,26]. GIRO ionograms are local/remote ionosonde measurements, not full end-to-end channel measurements [4]. KeepTrack's API is licensed CC BY-NC 4.0, requires attribution, and is restricted to non-commercial use; a deployment must review these terms and may substitute another permitted source [18]. Each evidence item retains source, timestamp, assumptions, age, coverage, and confidence.

The score remains an estimate against the installed service objective, not a
reachability verdict. Evidence sources retain separate semantics and
uncertainty; eligibility and local authorization remain independent checks.
The propagation-model comparisons summarized in Appendix~\ref{app:model-evidence}
show why calibration must match the intended path class, season, frequency,
and outcome metric.

The ranker can refresh its recommendations as new evidence arrives, but a new
score does not automatically mean a new radio setting. The link manager acts
only when an improvement is meaningful, the evidence is fresh, and policy
checks pass. It never retunes during PTT or an active transfer. With one
radio, Linkwise selects one channel for the next session; multiple concurrent
links would require separate radios, leases, and a different scheduler.



\section{Operating concept}

\subsection{Link opportunity set}
The system distinguishes two related outcomes:
\begin{enumerate}
  \item \textbf{Opportunity coverage:} how many eligible gateways/channels are currently known and assessed for the station.
  \item \textbf{Next-session utility:} the estimated value of the best single gateway for a particular queued message or installed service objective.
\end{enumerate}
The first is a portfolio-maintenance objective; the second is an execution decision. The system reports both. It must not describe a candidate as reachable merely because it is listed in RMS or predicted by VOACAP.

Winlink counts must be labeled rather than collapsed into a single “number of links.” Linkwise separately reports unique RMS gateway stations, distinct gateway/protocol/channel records, station-profile-eligible channels, recently reported gateways, attempted RF sessions, modem-connected sessions, and client-confirmed message transfers. A single gateway may contribute multiple channels, and a single radio can use only one RF channel at a time. The public Winlink status page describes report-age categories (normal, more than two hours, more than four hours, and removal after six hours) and explicitly cautions that its traffic display is not an exact message count [3]. These are source-specific status labels—not live RF reachability tests. The AI should always receive the inventory snapshot ID, source timestamp, query filters, metric window, denominator, and stale/coverage warnings alongside each count.

\subsection{Closed-loop operation}
Linkwise does not treat a directory entry or a forecast as a successful link.
It keeps the whole process visible—from observing candidate channels through
recording what happened after an attempt:
\begin{enumerate}
  \item Observe radio, modem, client, gateway, and propagation state.
  \item Validate source, freshness, and quality of each datum.
  \item Construct eligible candidate channels using station profile and service constraints.
  \item Rank routes for the objective, including uncertainty and switching cost.
  \item Select the best eligible route autonomously and prepare the radio under the installed authority.
  \item Read back and verify all configured state; confirm no active transfer/PTT.
  \item Have deterministic local policy automatically authorize an RF session when all configured checks pass.
  \item Start a bounded Winlink session, monitor it, and abort on safety or timeout.
  \item Record the result and update historical gateway estimates.
\end{enumerate}

\landscapeplantuml{build/diagrams/Linkwise_Operating_Loop}{Linkwise operating loop from evidence collection through session outcome and learning.}{fig:operating-loop}

\section{Safety, security, and reliability}

\subsection{Safety invariants}
\begin{enumerate}
  \item PTT key-up may be requested by the AI or modem, but can occur only through the arbiter under an automatically authorized session; key-down is always permitted.
  \item Policy is enforced in deterministic local code, not model instructions.
  \item Authorization is one-use, short-lived, and bound to channel/configuration/session limits.
  \item Frequency, mode, power, and service rules are checked locally before session start.
  \item No tune/switch while PTT is asserted or a transfer is active, except emergency abort.
  \item A watchdog, session timeout, autonomous abort, human-only emergency stop, and physical inhibit are independent of AI.
  \item Stale/unknown measurements are never converted to favorable defaults.
  \item Loss of AI or internet access triggers deterministic safe fallback/receive-only behavior and cannot prevent emergency stop.
\end{enumerate}

\subsection{Threats and failure modes}
Prompt injection, malformed tool arguments, stale RMS data, inaccurate propagation reports, external radio changes, multiple clients racing rigctld, modem/client crash, stuck PTT, unauthorized MCP access, and sensitive payload logging are in-scope threats. Controls include typed schemas, allowlists, capability-aware validation, single-writer leases, readback, freshness limits, watchdogs, loopback binding, audit trails, and payload-free default logging.

\subsection{Anti-oscillation}
Channel management uses configurable minimum dwell, switch hysteresis, switch penalties, retry limits, failure cooldowns, and deferral during active sessions. Numeric thresholds remain \todo{} until tested on real station hardware and realistic gateway sets.

The invariants above need observable acceptance evidence, not just design
intent. Section~\ref{sec:requirements-verification} states what should be
checked and how; TBD limits remain open until station-specific testing
supplies defensible values.


\section{Requirements and verification strategy}\label{sec:requirements-verification}

\clearpage
\begin{landscape}
\thispagestyle{plain}
\small
\begin{longtable}{@{}p{0.95in}p{0.55in}p{4.65in}p{1.0in}@{}}
\caption{Representative system requirements and proposed verification methods.}\label{tab:requirements}\\
\toprule
\textbf{ID} & \textbf{Prio.} & \textbf{Proposed requirement} & \textbf{Method}\\
\midrule
\endhead
OVR-F-001 & P0 & Report radio state with field-level support, freshness, and errors. & Test \\
OVR-F-002 & P0 & Route all radio writes through the policy-controlled path. & Inspect/Test \\
OVR-F-003 & P1 & Normalize RMS gateway/channel fields and source retrieval time. & Test \\
OVR-F-005 & P1 & Rank eligible gateways with evidence, uncertainty, and switching cost. & Analysis/Test \\
OVR-F-006 & P1 & Explain ranking factors, source provenance, freshness, and exclusions. & Test \\
OVR-F-007 & P0 & Enforce one exclusive radio-control lease. & Test \\
OVR-F-008 & P0 & Read back and verify configuration before marking a candidate prepared. & Test \\
OVR-F-009 & P0 & Prohibit retuning during PTT or active payload transfer except abort. & Test \\
OVR-F-011 & P0 & Automatically grant a bounded session when the installed authority and live interlocks pass; no session-by-session human approval. & Demo/Test \\
OVR-F-012 & P0 & Expire authorization at configured time/session/lease boundary. & Test \\
OVR-F-017 & P0 & Provide a human-only emergency stop independent of AI, MCP, network, and application availability. & Demo/Test \\
OVR-F-019 & P0 & Gate PTT key-up on active session, permit, lease, verified configuration, and local duration/duty limits; allow unconditional idempotent key-down. & Test \\
OVR-F-018 & P0 & Autonomously start, monitor, recover, and retry within preinstalled authority; no routine human intervention or controls. & Demo/Test \\
OVR-F-020 & P0 & On session/permit/lease/watchdog fault, request and audit PTT release and confirmation state. & Test \\
OVR-F-025 & P1 & Exercise production tool contracts using simulated adapters and virtual time. & Test \\
OVR-F-026--029 & P1 & Reproduce scenarios, isolate hidden truth, compare baselines, and report distinct outcomes. & Analysis/Test \\
OVR-F-030--031 & P2 & Qualify modem channel tests and inhibited hardware-in-loop behavior separately. & Test \\
OVR-Q-001 & P0 & Release software PTT and enter configured safe state within \todo{} after controller loss. & Test \\
OVR-Q-005 & P1 & Apply hysteresis, dwell, cooldown, and retry bounds to automatic selection. & Test \\
OVR-Q-009 & P1 & Claim modem compatibility only after bidirectional reference interoperability. & Demo/Test \\
\bottomrule
\end{longtable}
\end{landscape}
\clearpage

The full actor, use-case, and requirement traceability baseline is maintained in \url{docs/architecture/requirements-and-use-cases.md}. Early qualification should use a dummy rig backend and simulated adapters for stale data, concurrent controllers, readback mismatch, and PTT/session timeouts. Simulation evidence, modem channel qualification, and real hardware qualification are distinct; build checks are not interoperability evidence.


\section{Use cases}

\subsection{UC-1: Observe station and opportunity set}
The scheduler and agent obtain a snapshot. OLM reads radio, modem, client, and RMS state; attaches UTC timestamps and age; identifies unsupported fields; and returns a partial result if a subsystem is offline. No radio state is changed.

\subsection{UC-2: Rank gateway alternatives}
The installed station profile supplies message/service objectives and preferences. OLM filters unavailable, unauthorized, unsupported, stale, or out-of-profile channels; evaluates evidence per candidate; then gives the autonomous planner a ranking, confidence, rationale, and exclusions. When evidence is weak it reports that limitation rather than fabricating precision.

\subsection{UC-3: Prepare a channel}
The agent selects an eligible candidate. OLM acquires an exclusive lease, checks no active session/PTT, validates local limits, configures the full enabled CAT control profile through Hamlib, reads back supported values, and verifies modem/client readiness. Preparation cannot transmit. A mismatch invalidates the prepared state and triggers automatic recovery or a safe pause.

\subsection{UC-4: Authorize and conduct a bounded Winlink session}
Local policy automatically issues a one-session grant bounded by gateway, frequency, modem/client, power ceiling, queue scope, configuration hash, and expiry. No per-session operator approval is requested. The client begins the session; OLM monitors session/transfer/PTT and enforces time and duty ceilings. Completion, fault, expiry, loss of lease, emergency stop, or abort ends the attempt and records a summary.

\subsection{UC-5: Recover from failure}
On a failed connection, OLM releases PTT, classifies and records the failure, applies retry/cooldown, refreshes stale metadata, and re-ranks alternatives. It automatically prepares a new candidate and obtains a new policy grant if eligible; otherwise it waits or enters a latched safe state. It never reuses a previous grant.

\subsection{UC-6: Maintain link opportunity coverage}
At configured intervals, OLM refreshes RMS channels and external observations and recalculates candidate eligibility. A stable active channel is protected from score-driven churn; a material change is acted on at a safe session boundary without operator intervention.

\subsection{UC-7: Use propagation forecasts and evaluate reentry evidence}
At each planning cycle, the autonomous agent obtains a ranked candidate set. The ranker obtains VOACAP estimates for candidate path/frequency/time combinations, then combines them with fresh RMS availability, station profile, current reports, measured conditions, and empirical outcomes. The response identifies the likely planning window, input assumptions, score contributions, confidence, and stale/missing sources. Before preparation, the link manager refreshes state and obtains an automatic policy grant; forecasts do not themselves authorize transmission.

In parallel, a research service ingests versioned reentry predictions and associates them with independent ionosonde/sporadic-E observations and per-session results. The agent may see the experimental score and its uncertainty labeled as non-operational. A validation report compares it to the VOACAP/operational baseline using event and control windows, held-out data, and shuffled covariates. Until a pre-deployment evidence/validation gate is satisfied, reentry data cannot alter the live candidate order. Even after enablement, it is only a bounded ranking covariate, never a reachability claim or authorization input.

Together, these use cases define the exchanges the adapters must support. The
interfaces below keep those exchanges explicit and prevent the AI-facing API
from becoming a shortcut around radio policy.


\subsection{Goal-to-verification traceability}
Table~\ref{tab:goal-traceability} gives representative links from system
intent through requirements and use cases to the architecture elements and
evidence expected to verify them. The companion requirements baseline remains
the authoritative, complete traceability record.

\begin{table}[htbp]
\centering
\small
\caption{Representative goal-to-verification traceability.}
\label{tab:goal-traceability}
\begin{tabularx}{\linewidth}{@{}p{1.45in}p{2.2in}X@{}}
\toprule
\textbf{Goal} & \textbf{Trace links} & \textbf{Architecture owner and verification evidence} \\
\midrule
Fresh, explainable candidate choices & OVR-F-003, 005, 006; UC-1, UC-2, UC-6 & Inventory, evidence, and ranker; snapshot replay and checks of provenance, freshness, uncertainty, and exclusions \\
Bounded autonomous station operation & OVR-F-007--012, 018--020; OVR-Q-001; UC-3--UC-5 & Link manager, policy, radio adapter, and PTT arbiter; fault injection and inhibited hardware-in-loop demonstration \\
Measured session and transfer learning & OVR-F-025--031; OVR-Q-009; UC-4, UC-5, UC-7 & Modem, Winlink client, and history store; interoperability runs and denominator-bearing held-out evaluation \\
\bottomrule
\end{tabularx}
\end{table}

\section{Architecture}

The operating loop defines the behavior; the architecture assigns each step
to a component and sets boundaries between planning, policy, and radio control.
The views below use C4/PlantUML context and component diagrams alongside
deployment, data, and safety descriptions. They are complementary views of
one proposed system, not separate implementations.

\subsection{Architecture viewpoints}
The architecture is described from the concerns that matter to its users and
builders. These viewpoints organize the views that follow; the supporting
System Architecture Description contains the fuller stakeholder-to-view
mapping.
\begin{itemize}
  \item \textbf{Operational and functional:} how an operator and AI use the
        system, what its components do, and how a session moves from candidate
        selection to outcome recording.
  \item \textbf{Information and evidence:} how gateway, forecast, observation,
        radio, and session data retain source, time, uncertainty, and status.
  \item \textbf{Deployment and physical:} where the link manager, AI host,
        adapters, radio, and external services run, and how they connect.
  \item \textbf{Safety and security:} where trust boundaries lie and how
        authorization, PTT control, privacy, and recovery are enforced.
  \item \textbf{Development and verification:} how adapters and policy can be
        tested independently, including in simulation and controlled hardware
        tests.
\end{itemize}

\subsection{Context and components}
All architecture diagrams in this paper are authored as PlantUML source in
\path{docs/whitepaper/diagrams/} and rendered locally during the build. The
generated standalone LaTeX is an intermediate used to create vector PDF
figures; it is not maintained as diagram source [17].
\landscapeplantuml{build/diagrams/Linkwise_System_Context}{System context and component boundaries for the proposed Linkwise Link Manager.}{fig:system-context}

The propagation pipeline is an asynchronous application service, not work
performed inside an MCP request handler. Bounded source adapters collect
configured external and local evidence; normalization validates units,
timestamps, identities, and source-specific revisions before immutable
records enter the evidence store. The analysis service consumes a consistent
snapshot and produces versioned path analyses. The link manager combines
those analyses with current RMS inventory, station eligibility, and session
history to create ranked candidate links. MCP exposes these results to the
AI and accepts bounded refresh, preparation, and session requests. A candidate
is a proposed route for one eligible RMS channel; creating one never creates
a Winlink service, reserves spectrum, or authorizes transmission.

\landscapeplantuml{build/diagrams/Linkwise_Propagation_Evidence_Flow}{Proposed propagation evidence flow: independent collection, normalization, persistent analysis snapshots, candidate links, MCP queries, and policy-mediated sessions.}{fig:propagation-evidence-flow}

The current Rust radio-control MCP prototype is a limited implementation slice; its verified scope and remaining gaps are recorded in the lab-notes chronicle. It does not implement the complete autonomous link manager described here.

\subsection{Layered software design}
ChatGPT Agent or Codex supplies the external decision-making orchestration.
The local link coordinator executes bounded requests, owns authoritative
station state, and enforces timing and policy; it is not a competing strategic
planner. The detailed modem and DSP views in Sections~\ref{sec:modem-software}
and~\ref{sec:modem-dsp} refine this allocation.

The proposed software is organized as a ports-and-adapters system with a
deliberately controlled—not merely read-only—actuation boundary. The AI host
can discover typed operations and invoke the full locally enabled CAT and PTT
tool surface, while it has no direct device socket, shell, or route around
Linkwise policy. The MCP adapter validates request shape and identity, then delegates to application
services. Application services orchestrate transactions and deadlines; the
domain core owns eligibility, ranking, authorization rules, leases, and state
transitions. Outbound ports express capabilities without binding the domain
to Hamlib, a modem, a Winlink client, a database, or an LLM vendor. Adapters
translate those contracts to concrete protocols and report unsupported
capabilities explicitly.

This dependency direction is also a test strategy. Domain decisions can be
replayed against deterministic snapshots, and infrastructure can be exercised
through contract tests without a model in the loop. The simulator implements
the same ports as the station adapters but keeps evaluator-only RF truth out
of production-facing tools. The full responsibility and dependency allocation
is recorded in the companion Software Design Description (SDD) [20].

\landscapeplantuml{build/diagrams/Linkwise_Software_Layers}{Proposed layered software design, trust boundary, and allowed dependency direction.}{fig:software-layers}

\subsection{Component contracts and ownership}
The system uses typed commands and results rather than passing arbitrary JSON
maps into control logic. Requests carry a request and correlation identifier,
an idempotency key, an authenticated local principal, a deadline, and the
policy revision under which they were issued. Observations carry a source,
observation and retrieval times, validity, provenance, quality class, and
warnings. Missing capability or state is represented as unknown/unsupported,
not silently converted to zero or false.

The principal contracts are deliberately explicit. Radio control discovers
and exposes every backend-advertised field and supported operation, with
registered, schema-defined extensions for device-specific CAT features; writes
are serialized under a lease and returned with readback. PTT has separate
observation, assert, and release operations: the AI can request key-up, but
the arbiter validates its authorized session or test context and local limit;
release is idempotent and available during faults. A separate RF-environment
port ingests compact observations from a configured receive-only SDR/DSP
pipeline; it does not route IQ through MCP or control the IC-7200. The modem
reports lifecycle and link metrics; the Winlink client owns its supported
session and transfer operations. Neither modem nor client events alone prove
recipient delivery. Evidence providers return forecasts, reports, or local
RF measurements with provenance, never permission. Repositories persist
append-only decisions and lifecycle events; message bodies and credentials
are excluded by default. These contracts are summarized in Figure~\ref{fig:software-contracts};
the companion SDD defines their semantics and representative Rust signatures.

\landscapeplantuml{build/diagrams/Linkwise_Software_Contracts}{Proposed domain services and typed ports for CAT, policy-mediated PTT, modem, Winlink, evidence, and persistence.}{fig:software-contracts}

\landscapeplantuml{build/diagrams/Linkwise_Evidence_Model}{Proposed evidence-to-analysis model: collection runs, typed immutable observations, forecast inputs, inventory, station-profile assumptions, and historical outcome features.}{fig:evidence-model}

\landscapeplantuml{build/diagrams/Linkwise_Decision_History_Model}{Proposed decision and outcome lineage: individual ranked candidates, optional selection, policy evaluations, retries, and session lifecycle events.}{fig:decision-history-model}

\subsection{Command, state, and concurrency model}
Reads operate on immutable inventory/evidence snapshots and may run
concurrently. Station mutations are serialized by one station command actor
that owns an exclusive radio lease and monotonically increasing state revision.
This prevents a stale recommendation from overwriting a newer radio state or
autonomous command revision.
Preparation is a bounded saga, not an atomic radio transaction: Linkwise
revalidates the candidate and profile, checks idle/PTT conditions, applies
typed settings in a known order, and verifies supported readbacks. A
short-lived prepared handle is issued only for the verified state and retains
the lease until session start adopts it or expiry/cancellation releases it.
Partial failure triggers bounded compensation; unverified restoration faults
the station and blocks a new session.

Session start rechecks the candidate, profile, authorization, lease, radio
readback, inventory age, client/modem readiness, and interlocks. It does not
trust an earlier ranking or prepare response as authorization. A local policy
grant is issued automatically by deterministic policy inside the preinstalled
autonomous envelope; the model cannot mint or widen it. During a session, the
Winlink client and modem sequence protocol activity. Every modem-originated
keying request passes through the PTT arbiter, which checks session/config
binding, one-use permit, airtime limits, and watchdog state. Abort, timeout,
fault, lease loss, and permit expiry prioritize release. The sequence and
state views make these independent checks explicit.

\landscapeplantuml{build/diagrams/Linkwise_Session_Sequence}{Agent-orchestrated preparation and local deterministic authorization. Prepared handles bind verified configuration and expiry; the grant is a separate bounded policy result.}{fig:session-sequence}
\landscapeplantuml{build/diagrams/Linkwise_Agent_Session_Start}{Asynchronous session start and observation. Accepted work runs locally outside an individual MCP request, while the agent follows authoritative state through durable operation IDs.}{fig:agent-session-start}

\landscapeplantuml{build/diagrams/Linkwise_Session_State}{Proposed session lifecycle and fault transitions; release and abort remain available independently of normal progress.}{fig:session-state}

\subsection{Execution isolation and persistence}
Four execution responsibilities keep slow or unreliable work from delaying
safety. A safety supervisor owns permit expiry, key-up ceilings, session
deadlines, abort, and release confirmation. A station actor serializes
mutations. Bounded observation workers refresh directories and evidence
independently. An event writer persists decisions and lifecycle changes with
correlation IDs, source times, and adapter/policy versions. MCP workers are
bounded and rate-limited; safety actions use a priority path. A disk or audit
failure is surfaced rather than silently acknowledged, while physical PTT
release proceeds even if logging fails and raises a local alarm.

The local state machine and live adapter readbacks—not the model conversation
or a remote cache—are authoritative for an active session. On restart, all
prepared handles, grants, and permits are invalid. Linkwise reconciles PTT,
radio configuration, modem/client activity, and incomplete events before
enabling control, defaulting to receive-only/blocked until the state is known
safe. SQLite/WAL is a proposed initial store behind repository interfaces;
it is not a domain dependency. Counters are derived from deduplicated events
and retain their metric window, denominator, coverage, and source snapshot.

The persistence model separates immutable evidence from replaceable current
views. A collection run records provider/adapter version, completion status,
errors, and coverage. Each source record includes a typed measurement or
reported value as well as provider identity/revision, observation/issue and
retrieval times, validity interval, units, location/path/frequency where
applicable, quality, attribution, and a content checksum or permitted raw
record reference. Corrected provider records link to the version they
supersede. A forecast run records model/provider version, issue time, horizon,
complete input assumptions, typed outputs, and source-record references.
An analysis snapshot references the exact station-profile revision, inventory
snapshot, and evidence inputs used; each input identifies its feature role,
transformation, freshness, quality, and value. The snapshot also records
path/channel/time window, feature values, score components, uncertainty,
coverage, and analyzer/ranker/calibration versions.

A candidate ranking contains explicit ordered entries, one for each candidate
link, with rank, score components, eligibility, exclusion reasons, and the
analysis snapshot supporting that entry. A decision records whether the agent
selected a route, chose to wait, or found no eligible route; selection is
optional and references a specific ranking entry. Decisions and policy
evaluations are retained even when no session starts. One decision may lead
to multiple session attempts, including bounded retries. Each attempt keeps
connection result, failure category, transfer stage, delivery status when
confirmed, radio revision/configuration hash, and lifecycle events separate.
Historical aggregates reference the attempts used and record their matching
query, time window, sample count, numerator/denominator, estimate, and
uncertainty. Revisions create new records; they do not rewrite the evidence
used by a prior decision.

Index records by time, provider, path, gateway/channel, and session; retain
stable IDs and schema versions so exports can be replayed. Source-specific
licensing and terms determine whether raw payloads may be retained; where they
may not, retain only permitted normalized fields, attribution, hashes, and
references. Keep message bodies and credentials out of the evidence store.
MCP history reads return bounded summaries with counts, denominators, time
windows, coverage, and evidence IDs rather than unbounded raw history.
Retention, backup, migration, and export policies are deployment
configuration and must not silently discard records needed to explain an
active or past decision.

\subsection{Component responsibilities}
\begin{description}[leftmargin=1.4in,style=nextline]
  \item[AI planner] Interprets goals, compares explanations, and may invoke the full locally enabled CAT/PTT tool surface or request link workflows. It has no direct device connection and cannot alter local policy or self-approve key-up.
  \item[MCP interface] Discovers tools/resources and translates requests to the internal command API. Read-only, prepare, session, and abort actions are separate permission classes.
  \item[Policy engine] Validates profile, authorization, mode/frequency/power, channel freshness, lease, and session budget. It is deterministic and independent of model availability.
  \item[PTT arbiter] Owns the only software key-up path, routes authorized requests to qualified radio/modem PTT drivers, and enforces per-key-up/session budgets and watchdog release.
  \item[Link manager] Owns state machine, polling cadence, candidate lifecycle, hysteresis, radio lease, and session supervision.
  \item[Evidence collectors and normalizer] Poll configured sources on bounded schedules; validate, deduplicate, timestamp, and persist evidence without changing radio state.
  \item[Path analyzer and gateway ranker] Consume pinned evidence and inventory snapshots; compute versioned path analyses, eligible candidate rankings, score explanations, freshness, coverage, and uncertainty.
  \item[Winlink inventory/history services] Synchronize and deduplicate station/channel inventories, maintain source-age summaries, and aggregate session/transfer events without exposing message content.
  \item[Hamlib/PTT adapters] Expose the full capability-reported radio control set through typed calls (plus registered device extensions where needed) and a qualified PTT route (radio backend/CAT, serial control line, or modem/client request). The agent may request supported controls; all writes are serialized and audited, while every key-up is routed through the PTT arbiter.
  \item[Modem adapter] Manages state and the session interface for VARA HF or a future Linkwise modem. It does not own Winlink message routing.
  \item[Winlink adapter] Uses a documented client/list interface to refresh/use RMS channels, read optional aggregate queue state, and conduct sessions. Unsupported UI automation is not treated as a stable API.
  \item[Audit/state store] Persists station snapshots, evidence provenance, policy decisions, session summaries, and model correlation IDs. Message payload retention is off by default.
\end{description}

\subsection{Deployment}
The baseline installation is one operator workstation running the AI host, OLM, radio/modem/client adapters, and local persistence. The radio attaches via CAT/USB plus audio/PTT. Internet is used for RMS and optional propagation evidence, but PTT watchdog, release, emergency stop, radio lease, and transmit timeout are local. MCP, rigctld, and modem command ports are loopback-bound by default. Remote deployment requires authenticated transport and explicit security review.

\subsection{Evidence model}
Every item is tagged as \evidence{DOCUMENTED}, \evidence{MEASURED},
\evidence{VALIDATED}, \evidence{HYPOTHESIS}, or \evidence{UNKNOWN}. A
gateway status, a monthly path forecast, a spot report, an ionogram, and a
completed VARA session answer different questions. The ranker keeps those
distinctions visible rather than blending them into one unlabeled probability.
Every derived analysis is reproducible from immutable source and inventory
snapshot IDs plus the versioned model, ranker, and station-profile
assumptions. Historical summaries are computed by the link-history service,
not inferred from the model's conversation memory. Similar-session results
include the selection window, matching dimensions, numerator, denominator,
and uncertainty; sparse history is returned as sparse rather than promoted
to a confident success estimate. The AI may compare these summaries and
explanations, but deterministic application code owns canonical statistics
and ranking.

\subsection{Design choices and consequences}
Several choices shape the system and make its boundaries explicit. They remain
proposed until implementation and review confirm them.
\begin{enumerate}
  \item \textbf{Give the AI complete mediated device control.} Limiting the
        agent to recommendations would prevent it from managing station
        configuration and PTT end-to-end. Linkwise therefore exposes the
        backend's complete advertised CAT capabilities and explicit PTT
        requests through typed tools. A single local command path applies
        leases, operating profiles, authorization, auditing, and release
        watchdogs; the model cannot bypass these controls with raw device
        access or self-issued transmit permits.
  \item \textbf{Separate radio, modem, and Winlink adapters.} A single
        all-purpose adapter would be simpler initially, but each subsystem
        has different interfaces and failure modes. Separate adapters can be
        qualified independently, with the session controller coordinating
        them.
  \item \textbf{Optimize one active HF channel first.} Concurrent scheduling
        was considered, but one transceiver serves one channel at a time. The
        initial scheduler can therefore be evaluated against a realistic
        resource model; multi-radio scheduling remains future work.
  \item \textbf{Keep reentry features in shadow mode.} Adding them directly to
        the live score is premature because their predictive value is not
        established. Shadow evaluation tests the hypothesis without changing
        operational decisions.
\end{enumerate}

The views must agree on key allocations: the policy engine and session
controller run locally; every radio write and PTT key-up passes through their
controlled path; the AI host communicates only through MCP; and evidence
sources can inform ranking but cannot authorize transmission. These
correspondences are traced in the companion System Architecture Description.

\section{Simulation and decision evaluation}
Agent link-management decisions require repeatable evaluation against an RF
environment whose actual state is hidden from the agent. The proposed
simulation uses a fidelity ladder rather than asking one model to answer every
question. A fast Rust discrete-event harness first exercises the production
tool contracts with simulated radio, modem, client, gateway, and evidence
adapters under virtual time. Its evaluator owns latent channel quality,
correlated openings and fades, outages, interference, and eventual connection
and transfer outcomes; the agent sees only timestamped, potentially stale or
noisy observations. Simulator administration and truth inspection are test
runner operations, never production MCP tools.

Statistical path planning and modem testing serve different purposes.
ITU-R P.533 circuit predictions and local VOACAPL runs can inform priors and
baselines, but monthly/hourly reliability is not minute-by-minute ground truth
[12,15]. HF modem-in-loop testing should separately use documented
ionospheric channel cases such as ITU-R F.1487 [13]. GNU Radio's dynamic
channel model provides useful impairments, including fading, Doppler,
frequency and sample-rate offsets, and noise [14], but is not by itself proof
of a validated F.1487 implementation. The channel model, version, parameters,
and test vectors must be recorded. A later hardware-in-loop stage validates
actual CAT, audio, PTT readback, watchdog, and safe-stop behavior using an
inhibited, shielded or conducted radio path. No simulator stage is a license
or a reason to transmit over the air.

Decision experiments replay paired, seeded scenarios for the agent, simple
baselines (fixed ordering, round-robin, freshest/strongest evidence, and
historical success ranking), and a truth-informed oracle upper bound. Separate
random streams prevent policy actions from changing unrelated channel draws.
The oracle is a regret reference, not a deployable policy. A mission utility
values priority-weighted messages transferred by deadline while accounting
for airtime, retries, switching, and delay; safety rules remain hard
constraints, and any prohibited PTT event invalidates the run.

Reports should include eligible gateway/channel inventory, attempts,
connections, client-confirmed transfers, deadline delivery, time-to-transfer,
throughput, failed attempts, churn, airtime, queue age, calibration, regret,
tool errors, and safety outcomes. Include denominators, scenario coverage,
uncertainty intervals, and held-out scenario families; preserve traces and
manifests for replay. Calibrate channel transition and success distributions
against measured station history where available, preserving uncertainty when
data are sparse. Sporadic-E may be an explicit exploratory scenario overlay,
but neither it nor reentry timing is assumed to provide predictive power. To
test reentry-informed ranking, add paired ablation arms: VOACAP/operational
baseline; baseline plus reentry covariates; and shuffled-time/place reentry
controls. Evaluate only on held-out events, paths, and seasons, report
calibration and incremental mission utility with uncertainty intervals, and
keep the feature shadow-only unless it beats the baseline robustly. The
simulator must model uncertain prediction windows and event-source latency;
it must not leak hidden RF truth or future outcomes to the agent. Simulation
findings apply to the modeled scenario distribution and do not establish
real-world optimality. The detailed levels and proposed requirements are in
\texttt{docs/architecture/rf-agent-simulation-design.md}.
These limits matter: a successful simulation demonstrates performance under
its assumptions, not that the same result will hold on the air.

The architecture becomes easier to evaluate through the work the autonomous
agent must perform. The use cases follow its continuous loop: observe, compare,
configure, authorize locally, operate, recover, and learn. The human runtime
role is limited to the emergency stop.

\section{Interfaces and control contracts}

\subsection{Radio-control adapter: Hamlib/rigctld}
The radio adapter exposes the complete control catalog advertised by the selected Hamlib backend, subject to the installed station profile. For the target IC-7200, Hamlib/rigctld is the CAT boundary; supported operations and meter levels must be discovered from the installed backend at runtime rather than inferred from the model name [22, 45]. The adapter does not promise that every Hamlib operation is supported or readable on this radio.

The read-only tools are \path{radio_get_state} (optional maximum age),
\path{radio_get_capabilities}, \path{radio_get_controls} (optional
category/control filters), \path{radio_get_control} (exact catalog ID), and
\path{radio_get_ptt_state}. Results include radio identity, observed values,
per-field support/errors/freshness, control ranges/enums, readback status,
and backend capability provenance. Signal-strength levels, if exposed, retain
the Hamlib level token and radio-specific units/calibration; they are local
receiver measurements, not modem link evidence.

The mutating tools are \path{radio_set_control} and
\path{radio_invoke_extension}. The first accepts an exact discovered control
ID, a matching tagged value, and an optional expected state revision. The
second accepts a registered extension ID and schema-validated arguments;
arbitrary CAT bytes, shell commands, and direct rigctld requests are never
accepted. Writes acquire/validate the station lease, apply profile limits,
serialize CAT access, audit the request, and read the value back where
supported. The response distinguishes accepted request from verified state;
uncertain outcomes require a fresh observation. Access class and side-effect
class are explicit per control. The IC-7200 target is not yet physically
qualified, and the prototype's current local revision does not detect
changes made by an external controller.

\subsection{MCP shared contract}
The MCP API separates observation from actuation and maps each purpose-specific
tool to a typed Linkwise port. Results include request/correlation IDs,
status, observation time, data, warnings, and errors; evidence retains source,
source and retrieval times, units, quality, freshness, and provenance.
Unknown, stale, unsupported, unavailable, and zero are distinct. Collection
and history reads are bounded, and provider refreshes use preconfigured
sources only. Exact schemas, resource URIs, common errors, and versioning are
defined by the companion MCP interface design. MCP annotations are advisory;
local deterministic policy authorizes every side effect.

Each adapter section below defines its MCP surface, result evidence, and
capability limits. The link-history adapter correlates their observations by
session and attempt without conflating local RF activity, modem connection,
Winlink session completion, CMS acceptance, and recipient delivery.

\subsection{Actuation and link-session coordinator}
The mutating radio tools are \path{radio_set_control} and
\path{radio_invoke_extension}; their exact value schemas, side-effect classes,
lease checks, audit records, and readback behavior are described in the
radio-control adapter contract above. Link workflows use
\path{link_prepare_candidate}, \path{link_request_session_authorization},
and \path{link_start_session}. Preparation re-fetches the candidate,
acquires the station lease, confirms idle/PTT state, applies radio settings,
verifies readback, and returns a short-lived handle bound to that state.
Authorization is a separate local deterministic policy decision, bounded by
candidate, configuration, service scope, power, and expiry. Starting a
session consumes both handles; a prior MCP read is never authorization. MCP
discovery/invocation does not itself grant authority [5].

\subsection{PTT and transmit-arbitration adapter}
PTT is separate from CAT configuration because it is a safety-critical
transmit path. \path{radio_get_ptt_state} reports configured route,
requested/observed state, readback support, watchdog status, and observation
time. \path{radio_ptt_request} accepts \path{key_up} or \path{key_down};
key-up is bound to an active session and one-use permit, with no caller-
selected duration or duty limit. \path{link_request_ptt_permit} revalidates
session authorization, lease, radio configuration, interlocks, and remaining
budget before issuing that permit. Key-down is idempotent, always available,
and does not require a session or permit.

The modem and client request burst timing only through the same local
arbiter. It verifies the route and permit, enforces configured continuous-
key-up and total-session limits, confirms release when readback is available,
and records each transition. Watchdog, timeout, lease loss, abort, controller
loss, and emergency stop all cause release; inability to confirm release
produces a latched fault/inhibit according to station policy. The human
emergency stop is physically independent and is not an MCP tool. These are
proposed guarantees: the current Rust prototype denies key-up, has no permit
arbiter or watchdog, and is not transmit-qualified.

\subsection{Audio-device and routing adapter}
The audio adapter discovers and reports the host audio endpoints used by the
IC-7200 USB codec or a qualified external interface. Read-only
\path{audio_get_status} reports selected input/output routes, stream state,
sample format/rate, channel count, last buffer time, underrun/overrun counts,
and measured route health. \path{audio_get_capabilities} lists available
devices and supported formats without exposing OS handles or arbitrary
device paths. Status fields distinguish device-present, stream-open, and
audio-observed states; an open stream is not proof that audio is correctly
decoded by the modem.

Route selection is limited to preconfigured, qualified profile IDs. If
runtime selection is enabled, \path{audio_set_route} can select only a
profile while PTT is released and no session is active; it returns the
effective route and readback/stream validation. Arbitrary device names,
filesystem paths, gain values, and routing graphs are not MCP inputs. The
modem adapter consumes the route's readiness and health signals, while the
link coordinator owns when a route change is permitted. The current Rust prototype has optional CPAL input capture and bounded
receive workers, but no qualified playback/routing service or audio MCP tools.
Section~\ref{sec:modem-dsp} details the implemented input and DSP paths.

\subsection{RF-environment observation adapter}
This read-only adapter ingests compact results from a configured
receive-only SDR and DSP pipeline, such as GNU Radio/SoapySDR. It supplements
the IC-7200's tuned-channel observations; the IC-7200 USB audio is demodulated
audio, not a wideband IQ feed, so it cannot survey other channels while
serving an active link. The adapter does not tune or key the IC-7200 and does
not control arbitrary SDR parameters at runtime.

Its question is whether a configured receiver observed measurable RF
activity in a specified span and interval, and how trustworthy that
measurement is. It does not decide whether a peer decoded a frame or whether
a protocol link was established. Spectrum occupancy, calibrated noise floor,
energy bursts, and channel-power summaries require an IQ-capable receive
source and are reported only when supported by that source and profile.
Measurements derived from IC-7200 USB audio are audio-domain observations of
the tuned, demodulated channel; they are not interchangeable with RF power or
a wideband channel survey. Any uncalibrated audio-level or SNR-like estimate
must be labeled as a proxy with its estimator and limitations.

\path{rf_get_status} reports receiver/collector health, configured profiles,
last collection time, and coverage. \path{rf_get_capabilities} describes
supported bands/spans, metrics, resolution, sample configuration, and
calibration state. \path{rf_get_snapshot} reads the latest immutable
observation for a configured profile; \path{rf_get_history} returns a
bounded time range and result count. A separately authorized
\path{rf_request_observation} may request one bounded capture using a
preinstalled profile ID and dwell limit; it cannot supply arbitrary
frequencies, gain, sample rate, or device path. Results may include noise
floor, spectrum bins or summaries, occupied-bandwidth/energy candidates,
receiver gain/configuration, frequency span, units, calibration/quality
flags, sample interval, estimator identity, capture source, uncertainty, and
missing-coverage reasons. Capabilities identify which of these are available
for each profile; unavailable or invalid measurements are represented as
such, never as zero. Results include capture start/end times and freshness so
the consumer can distinguish a current observation from stale data. Raw IQ
and continuous waterfall frames remain local to the DSP pipeline and are
never MCP tool payloads or ordinary audit records. Energy detection indicates
local activity only; it does not identify a callsign, prove decoding, or
establish gateway reachability. During local transmission, desense, or a
receive gap, the adapter reports degraded or absent coverage rather than
interpreting that interval as a quiet channel.

\subsection{Modem adapter: VARA or future Linkwise modem}
The modem adapter owns waveform/modem lifecycle and protocol-level link
state, not Winlink message routing or radio CAT. \path{modem_get_state}
reports process/device readiness, audio/PTT route health, and current state.
The \path{linkwise://modem/state} resource exposes the latest read-only
snapshot with source time and freshness;
\path{modem_get_capabilities} reports only verified operations and fields
for the configured modem/version. \path{modem_get_session} returns one
session's state and negotiated parameters where available.
\path{modem_get_events} accepts a session ID, cursor, and bounded limit to
return ordered acquisition, peer-decode, handshake, link-up/down, quality,
and termination events. Event data preserves native event code, normalized
event class, peer identity only when actually decoded, source time, and
adapter version. A link-up event means the modem's protocol state machine has
completed the protocol's required bidirectional handshake and entered its
connected state. Carrier/energy detection, synchronization or acquisition,
one decoded peer frame, or a locally issued connect request is only evidence
of an earlier stage and must not be normalized as link-up. The session
snapshot records the protocol state, transition time, peer identity when
provided by the modem, negotiated parameters, and termination reason when
available. BER, SINAD, SNR, throughput, or retry counters are included
only when the selected modem actually exposes them, with native units and
measurement semantics; Linkwise does not infer missing values.

Session lifecycle commands are issued through the link coordinator, not a
raw modem-control MCP tool. The adapter may offer internal start/stop/status
operations to that coordinator over a verified client interface. VARA's
external command/status behavior must be qualified against the selected
release before its fields become stable MCP contracts. For VARA, Linkwise
must consume a qualified explicit status/event from the selected release and
preserve its native code and semantics; it must not infer connection from
audio, RF energy, or a locally requested connect. A future native modem uses
the same domain event contract, with its distinct implementation version.
The evidence layers are intentionally separate: RF activity is a DSP
observation; modem link-up is a protocol connection; Winlink session and
transfer completion come from the Winlink client; CMS acceptance or recipient
delivery is reported only when the client explicitly confirms it. Modem
link-up alone proves none of the latter outcomes.

\subsection{Winlink RMS inventory adapter}
This adapter refreshes and normalizes the configured published RMS channel
source. \path{winlink_refresh_inventory} accepts optional maximum age and
preconfigured protocol/service filters, and returns an immutable snapshot
ID, source/retrieval times, per-source status, warnings, and count deltas.
\path{winlink_get_inventory_summary} returns unique gateway/channel counts
and status/report-age buckets. \path{winlink_get_gateway_channels} accepts
a gateway identity and returns its listed channel frequencies, modes,
services, source identifiers, status, and freshness. The
\path{linkwise://winlink/channels} resource exposes the corresponding cached
channel view with snapshot identity and freshness; it does not perform a live
probe.

Inventory records identify callsign/station, channel frequency, protocol,
service, location/grid where published, source timestamp, local retrieval
time, and normalization warnings. Directory status and report age are
discovery evidence only. They do not establish that the gateway is currently
on air, that the HF path is open, or that the local station can reach it.
Refresh performs bounded network I/O against configured sources and cannot
accept arbitrary URLs, credentials, or polling schedules. A partial refresh
preserves the prior snapshot as stale and reports the failure rather than
replacing unavailable counts with zero.

\subsection{Winlink client and session adapter}
The client adapter owns queue/session/message-service interaction and keeps
those outcomes separate from modem link state. \path{winlink_get_status}
reports client readiness, supported connection methods, and modem/client
association. \path{winlink_get_queue_summary} returns only aggregate counts
or bytes for supported clients, with queue direction, observation time, and
capability status; payloads, recipients, and message contents are excluded
by default. \path{winlink_get_session} returns client session ID, selected
RMS gateway/channel, lifecycle stage, timestamps, transfer progress when
available, and terminal reason. \path{winlink_get_events} returns bounded
ordered session/transfer events. Deduplicated cross-session outcomes are
queried through \path{link_get_history}, which omits message contents by
default. The \path{linkwise://winlink/queue-summary} resource exposes only
the supported aggregate queue view; session resources remain addressed by
session ID.

Session start and abort are routed through the link coordinator and local
authorization and PTT arbiter, not directly from the inventory or modem
adapter. A client adapter is implemented only against a verified supported
interface; Pat remains a candidate, and the design does not assume Winlink
Express GUI automation. Outcomes distinguish gateway contact, RF/modem link,
client session completion, CMS acceptance, and recipient delivery only if
explicitly confirmed by the client.

\subsection{Propagation and gateway data}
\path{propagation_get_status} returns provider health, latest source and
retrieval times, coverage, and adapter errors.
\path{propagation_get_path_analysis} accepts a station/path, channel or
frequency, UTC planning window, and optional maximum age; it returns an
immutable analysis ID, exact evidence input IDs/roles, model/config and
station-profile revisions, estimates, uncertainty, coverage, and
stale/missing-input reasons. \path{propagation_refresh_sources} requests a
bounded refresh from preconfigured providers and returns collection-run and
evidence IDs with per-provider outcomes. It cannot change polling policy or
name arbitrary URLs. VOACAP forecasts, PSK Reporter/DX Maps spots, GIRO
measurements, and space-weather sources remain separately attributed
evidence classes; none acts as a reachability guarantee or transmit permit.
Scheduled collection and analysis run independently of MCP calls.

\subsection{Link coordination, history, and evidence adapter}
\path{link_rank_gateways} accepts a service objective and permitted
band/mode/filter set, then returns an immutable ranking snapshot with
gateway/channel IDs, eligible/excluded status, component scores, evidence
references, freshness, uncertainty, and reasons.
\path{link_get_candidate_analysis} returns the pinned analysis and
ranking-entry IDs, input provenance, and any associated policy decision.
\path{link_get_similar_sessions} accepts bounded matching dimensions and a
UTC window; it returns measured outcome counts, denominators, sample size,
coverage, and uncertainty, retaining sparse data as sparse.

\path{link_get_session}, \path{link_get_history}, and
\path{link_get_metrics} return ordered attempts and distinct connection and
transfer numerators/denominators with data coverage. Retries create new
attempts; session records are append-only and omit message payloads. The
correlated evidence view distinguishes radio configuration/readback,
transmission attempt, valid peer response, modem link establishment, Winlink
session completion, CMS acceptance, and message-transfer confirmation.
Energy detection, an S-meter value, PTT assertion, or RMS directory listing
cannot independently establish a link. The full machine-readable schemas,
resources, common errors, and versioning policy remain defined by the
companion MCP interface design; this section specifies each adapter's
architecture-level capability boundary.

These interfaces cross trust boundaries: the model can ask, external sources
can inform, but only local policy can authorize side effects. The safety and
security design makes that distinction enforceable.

% Detailed modem view. Included from the main whitepaper; not a standalone paper.
\section{Agent orchestration and modem software architecture}
\label{sec:modem-software}

This section connects the proposed ports-and-adapters architecture to the
Rust modem implementation as of 11 November 2026. The radio-free modem
preview serves traditional clients over paired TCP sockets and exchanges
audio with a bench peer. The proposed station runtime adds policy, device
ownership, and physical transmit control. The evidence table records
modem interoperability, client transfer, and station qualification separately.

\subsection{Standalone modem product and shared runtime}
\label{sec:modem-shared-runtime}
\emph{Proposed integration.} The native modem should be delivered as a shared
Rust runtime usable by traditional amateur-radio applications without an AI
host. The foreground executable \path{linkwise-modemd} will
assemble the runtime and its adapters. Optional operating-system supervision
and graphical configuration do not define a different protocol engine.
Linkwise orchestration remains an additional client of the station execution
contract, not a prerequisite for an ordinary client connection.

The common core owns session state, queues, deadlines, handshake and ARQ
sequencing, byte accounting, and outcomes. Platform adapters provide audio
capture and playback, device discovery, PTT, permissions, and process startup.
The implementation supports virtual time for deterministic tests and a
monotonic clock for the foreground runner. The device execution tests use
simulated audio and PTT.

The radio-free preview supports uncompressed BW500 and fixed level-4
BW2300/BW2750. BW500 remains the first target for station integration.
A release manifest must enumerate the qualified roles, DATA forms,
host commands, reply/query forms, rate behavior, and disconnect behavior.

\subsection{Traditional-client interface and session admission}
\label{sec:modem-client-admission}
\emph{Proposed integration.} A VARA-compatible host adapter exposes a
CR-delimited command/notification TCP channel and a separate raw binary DATA
channel. Ports default to loopback 8300/8301 where available and remain
configurable. These are host-protocol compatibility goals; neither accepting
a command nor opening a socket establishes on-air interoperability. Winlink
message framing, mailbox exchange, and message acceptance remain the client's
responsibility. No message body or continuous modem stream traverses MCP.

The preview uses \path{host::tcp} to pair one host command socket and one
DATA socket. \path{api::CommandDecoder} parses the command stream.
A separate IPv4 loopback socket exchanges audio with a single cable peer.
Each record has a four-byte big-endian sample count followed by little-endian
32-bit float samples at 48 kHz mono. The audio bridge can attach named
virtual capture and playback devices.
A compatibility matrix must record initialization commands, replies,
asynchronous notifications, buffer semantics, and unsupported behavior for
each qualified client/version. The current parser supports only a subset;
unsupported commands continue to fail explicitly. In particular, an
\texttt{OK} reply acknowledges command acceptance, not a completed RF link.

One station execution service arbitrates both admission paths:
\begin{description}
\item[Standalone.] A locally configured client principal requests
\texttt{CONNECT}. The service validates identity, station configuration,
route readiness, lease availability, and installed authority, then creates
one bounded operation and native session. Listening also requires an installed
receive/respond policy; a decoded inbound call does not itself authorize PTT.
\item[Managed.] The coordinator prepares and authorizes an attempt and
reserves an expected client association. The accepted start arms that
association; the matching client \texttt{CONNECT} consumes it once and starts
the backend. It must not create a second operation. The reservation binds
client ownership, local/remote callsigns, configuration revision, lease, and
expiry. Mismatches, competing requests, and attachment timeout fail explicitly.
\end{description}

The legacy two-socket protocol supplies no authenticated pairing token.
Initially permit one exclusive control/data pair per configured endpoint,
with a bounded attachment timeout and an explicit pairing state machine that
handles either arrival order. RF admission requires a complete pair.
A replacement socket cannot adopt an existing
session; when either channel closes, both are detached, admission is blocked
until cleanup completes, and the active session is cancelled with PTT release
requested. Peer address alone is not client identity. Loopback is a deployment
restriction rather than authentication; shared or remote deployments require
an authenticated enclosing transport or a separately qualified ownership
mechanism. Model disconnect is different from modem-client channel loss.

\subsection{Native session and device execution contracts}
\label{sec:modem-execution-contracts}
\emph{Proposed integration.} One serialized native session actor owns the
current peer, protocol phase, transmit/receive epoch, outstanding DATA, and
monotonic deadlines. Its lifecycle distinguishes idle, admitting,
establishing, connected, closing, completed, aborted, and faulted states;
transmit, answer-wait, query, and turn phases are subordinate protocol states.
Only validated events in the current phase may advance it. Every asynchronous
audio or device result carries a session identity and a unique transmission
token; results from a previous session or transmission are ignored and
recorded. Stop and channel-loss cleanup are idempotent. An accepted abort
invalidates outstanding work before any late completion can advance state.

Application lifecycle requests enter through \path{ModemSessionPort}; host
commands enter through \path{host::ModemPort} and the same admission service.
The native backend adapter binds both to the one actor. It must not run a
second protocol machine in the TCP handler or in the station coordinator.
External VARA and simulation remain alternative backends behind the lifecycle
port; only the native backend assembles the Rust DSP/session components.

Audio contracts identify a continuous capture epoch, absolute sample ranges,
qualified route and format, and terminal discontinuity. Playback accepts
owned bounded buffers and reports accepted, started, presented, cancelled,
or failed with the transmission token. Queue acceptance or callback buffer
consumption is not presentation completion. Cancellation must establish
whether queued samples can still reach the device; uncertainty faults the
session and prevents a new transmission. The PTT contract separates permit,
assertion, release request, and observed release. Actual presentation plus
qualified release/turnaround evidence is required before the actor emits a
protocol transmit-complete event and arms its receive deadline. Independent
watchdog release does not wait for the actor, event storage, or an AI host.

All mailbox, audio, event, and byte queues have configured finite limits.
Backpressure suspends acceptance at the host boundary without blocking DSP
callbacks or the actor's stop/watchdog processing. Protocol deadlines cannot
be extended by a slow client or a full notification queue. A stalled consumer
has a bounded timeout and a terminal reason; events needed to account for
bytes cannot silently disappear. Device loss, capture gaps, lease expiry,
playback timeout, and unconfirmed PTT release terminate the normal path.

Byte accounting distinguishes cumulative accepted bytes, pending bytes,
unique outstanding bytes, peer-acknowledged bytes, unresolved bytes, and
explicitly discarded bytes. For an operation, accepted bytes equal pending
plus outstanding plus acknowledged plus unresolved plus discarded bytes;
a byte occupies only one of those categories at a snapshot. Transmit attempts
and retries are separate counters and never inflate acknowledged totals.
Receive accounting distinguishes bytes decoded, reassembled/deduplicated,
queued for the host, and written to its socket. Socket write completion is
not proof the client consumed or committed the bytes. Host \texttt{BUFFER}
notifications need a client-qualified mapping to these counters rather than
the current ARQ queue size. Peer acknowledgment is not Winlink message or
CMS acceptance.

On ambiguous failure the operation result records unresolved byte counts and
ranges; payload remains in bounded session-owned memory rather than audit
events. Explicit abort/clear can discard bytes with a recorded count. No reconnect
or process restart replays them automatically. In-memory retention does not
promise crash durability: a crash produces an unknown outcome unless a
separately implemented bounded journal supplies evidence. Native on-air
duplicate handling must be qualified independently of local callback tokens.

\subsection{Shared-core implementation order and exit criteria}
\label{sec:modem-core-plan}
The first two stages have an implementation at the validated-event boundary
with a clock-driven runner and simulated devices. The BW500 audio loopback
and paired TCP bench exercise stages three and four with simulated peers.
Rolling pilot acquisition for levels 1--3 has separate off-air receive
evidence (Table~\ref{tab:modem-evidence-gates}). The remaining work includes
production DATA sequencing and recovery, native audio/PTT integration,
and broader client and reference-peer qualification. The sequence below defines the
scope and exit criteria for each stage.
\begin{enumerate}
\item \textbf{Session contracts and owner.} Define typed commands/events,
state transitions, operation/session/transmission identities, finite budgets,
byte-accounting rules, and a virtual clock. Demonstrate rejection of competing
owners, illegal transitions, stale completion, and repeated cleanup.
\item \textbf{Simulated audio and PTT execution.} Bind capture epochs,
playback tokens, presentation, release, and turnaround to the actor. Inject
late/missing completion, capture gaps, cancellation races, device loss, and
watchdog expiry. Demonstrate independent release and no false TX completion.
This yields the first foreground shared-core milestone without native devices.
\item \textbf{First complete link profile.} Integrate establishment,
packetization, DATA fields, validated replies/queries, ARQ, duplicate handling,
turn/rate ownership, keepalive, and disconnect. Exchange arbitrary binary
streams between two runtimes under deterministic loss, duplication, and delay;
preserve byte-accounting invariants and identify unresolved outcomes. Resolve
wire uncertainties with independent evidence rather than invented sequencing.
\item \textbf{Traditional-client service.} Add loopback listeners, exclusive
pairing, admission, notifications, timeouts, and bounded backpressure. A host
harness must exercise initialization, both socket arrival orders, competing
clients, binary preservation, stalled consumers, channel loss, and reconnect.
Only advertise the tested host-protocol subset.
\item \textbf{Station integration.} Bind standalone and managed admission to
policy, leases, grants, reservations, and operation/event records. Prove a
matching managed \texttt{CONNECT} starts exactly one attempt and a mismatch
cannot transmit. Test client loss separately from model loss.
\item \textbf{Reference and client qualification.} Demonstrate bidirectional
transfers with an unmodified release-specific reference modem and a named
Winlink client/version. Preserve commands, both audio directions, timing,
payload comparisons, recovery evidence, and distinct modem/client outcomes.
Native hardware measurements still require the corresponding device adapter.
\end{enumerate}

Only after these shared contracts stabilize should platform work supply
macOS audio/permissions and launch-agent integration, Linux audio/device
permissions and user-service integration, and Windows audio/device access
and background-process packaging. Foreground operation comes first on each
platform. A system service is optional and requires qualification in its
actual execution context; it is not evidence of audio readiness. Each native
adapter must satisfy the same fault/completion contract and undergo device
and inhibited/conducted PTT qualification before operational RF use.

\subsection{ChatGPT Agent and Codex as the orchestrating agents}
ChatGPT Agent or Codex is the decision-making orchestrator: it observes station
readiness, compares eligible RMS channels and evidence, chooses the next
attempt, invokes MCP tools, and adapts its plan to recorded outcomes. These
are alternative AI-host clients of the same typed application contract;
client transport and reconnect behavior must be qualified for each deployed
host. Neither client is a modem timing source, a radio driver, or an authority
issuer. The agent's context and conversation history are not station state.

The local link coordinator is the \emph{station execution service}. It performs
bounded use cases under deterministic policy rather than operating a second,
competing strategic planner. The agent may choose a candidate; the domain
ranker supplies reproducible scores and exclusions; local policy decides
eligibility and authorization; the execution service serializes preparation
and supervises the accepted session. Handshake, ACK/NAK, retransmission,
audio-buffer, and PTT deadlines stay local even while the agent is thinking,
disconnected, or recovering context.

The preparation contract remains the one in the interface section:
\path{link_prepare_candidate} returns a short-lived handle bound to candidate,
lease, verified configuration revision, route profile, and expiry;
\path{link_request_session_authorization} asks deterministic policy for a
separate bounded grant; \path{link_start_session} consumes both after a fresh
validation. Authorization is automatic within installed station authority,
not a runtime human approval step. Every mutation carries principal,
request/correlation IDs, policy revision, deadline, and an idempotency key.
Repeating an accepted start request must return the same operation rather than
create another RF attempt.

An accepted start returns a durable operation ID and session identity promptly.
The agent follows bounded session snapshots and ordered cursor-based events;
it does not hold one MCP request open for an entire contact. Durable operation
records, terminal reasons, and explicit unknown/stale fields let a reconnecting
agent reconstruct what happened. MCP disconnect and application restart are
different events: local policy determines whether an already accepted session
finishes or aborts after client loss, while process restart invalidates
in-memory grants/leases and requires reconciliation. Neither event extends a
transmit permit or authorizes automatic replay of unresolved bytes.

\subsection{Ports, ownership, and backend interchangeability}
The inbound MCP adapter validates identity and request shape and translates
results. Application services own preparation, deadlines, operation records,
and lifecycle orchestration. Domain services own policy, leases, session
invariants, and authorization. Concrete adapters own device/protocol details.
The proposed modem backend contract is \path{ModemSessionPort}; the existing
\path{host::ModemPort} is an inbound host-command/data boundary and must not be
mistaken for this outbound driver interface.

\begin{longtable}{@{}p{0.25\textwidth}p{0.69\textwidth}@{}}
\caption{Ownership at the agent, station, and modem boundaries.}\label{tab:modem-ownership}\\
\toprule
\textbf{Owner} & \textbf{Responsibility and boundary}\\
\midrule
\endhead
AI host & Select objectives and subsequent attempts from observed evidence;
request typed actions without minting grants, changing installed limits, or
addressing hardware sockets.\\
Station execution service & Own attempt/session identity, prepared handles,
lease adoption, client/modem association, cancellation, reconciliation, and
terminal operation results.\\
Modem backend adapter & Report capabilities and readiness; accept internal
bounded lifecycle requests; expose authoritative protocol events and a
backpressured data endpoint. Implementations may be native Rust, qualified
external VARA, or simulation.\\
Native session owner & Own one protocol state machine, receive/transmit epoch,
current peer, timers, outstanding bytes, and the ordering of DSP and playback
work. The simulated runtime implements session ownership and deadlines.\\
Audio adapter & Own configured input/output routes, formats, buffers, sample
continuity, and presentation-completion evidence. CPAL provides input capture and an output
primitive with a presentation estimate. Hardware completion remains open.\\
PTT arbiter & Own the only software key-up path, one-use permits, route checks,
key-up and session budgets, watchdog release, and requested/observed state.\\
Winlink client adapter & Own its supported queue, message-service session,
transfer semantics, and confirmation evidence. It does not delegate message
framing or body bytes to the language model.\\
\bottomrule
\end{longtable}

\landscapeplantuml{build/diagrams/Linkwise_Modem_Runtime}{Proposed agent-to-modem ownership and backend choices, over the native components currently implemented. Dotted integration edges remain incomplete; this is not an as-built operational runtime.}{fig:modem-runtime}

The backend port should provide versioned capabilities, an observed state
snapshot, internal begin/stop operations, a session/event cursor, and a bounded
data-stream attachment. The native adapter and external VARA adapter must
produce the same domain event classes while preserving native event codes,
backend versions, units, timestamps, unsupported fields, and qualification
provenance. A simulator implements that same port without exposing hidden RF
truth through MCP. Backend replacement cannot silently increase station
privileges or normalize an unverified signal into a connection.

\subsection{MCP observations and the modem data plane}
The modem's MCP surface is observational:
\path{modem_get_capabilities}, \path{modem_get_state},
\path{modem_get_session}, and \path{modem_get_events}, plus the cached
\path{linkwise://modem/state} resource. Session start/abort is routed through
the coordinator, and explicit CAT/PTT tools remain available through their
policy-mediated contracts. Audio-route selection uses qualified profile IDs,
never arbitrary device names, filesystem paths, gain scripts, or routing
expressions supplied by the agent.

Capability results must distinguish implementation from qualification: for
example, level-4 single-frame synthesis exists, but stock-peer acceptance and
multi-frame operation are unqualified. An advertised capability needs a
backend/version, mode, direction, supported parameters, availability, and a
qualification scope. State should distinguish device present, stream open,
audio observed, synchronized candidate, CRC-valid frame, handshake complete,
and active data exchange. Observations include sample ranges where useful and
source/freshness metadata; missing BER, SNR, throughput, or retry measurements
remain unsupported, not zero.

Protocol bytes travel between the Winlink client and modem on a bounded
stream or spool. MCP carries aggregate queue counts, message references where
supported, progress, and outcome summaries; raw audio, continuous IQ, and
message bodies are excluded by default. The agent-facing operation ID, station
attempt ID, modem-session ID, and Winlink-session ID remain distinct and are
correlated in the event store. UTC timestamps support audit correlation;
monotonic instants enforce deadlines; absolute sample counters locate signal
evidence. One must not substitute for another.

\subsection{Native implementation map and current limits}
The Rust package contains the foreground bench service and reusable components.
\path{audio::device} captures input; \path{modem::receiver} bounds and dispatches
receive candidates; \path{modem::waveform} supplies energy detection and tone
correlation; \path{modem::bw500} implements index-modulated levels 1--3;
\path{modem::bw500_level4} handles level-4 channel coding;
\path{modem::bw500_waveform} adds measured-reference two-band synthesis and
reception; \path{modem::wide_level4} supplies level-4 setup and level-2 through
level-4 DATA waveforms for BW2300 and BW2750, using the carrier allocations
in \path{modem::wide_allocations}. \path{modem::coding} implements whitening, CRC, interleavers, and
FEC. \path{modem::control}, \path{modem::setup}, \path{modem::handshake}, and
\path{modem::audio_handshake} cover keyed control, caller setup, and
establishment. \path{modem::data} and \path{modem::arq} are frame/container
and sender components. \path{api::CommandDecoder} and \path{host} implement
host stream boundaries.

The clock-driven \path{modem::runtime} combines establishment and ARQ at
the validated-event boundary. The virtual host in \path{linkwise-modemd}
owns its bench handshake, transmit turn, pending DATA, and deadlines.
The host keeps command and binary DATA streams on paired TCP sockets. A
separate virtual audio connection carries a big-endian sample count followed
by little-endian 32-bit floating-point samples at 48 kHz. The host queues
up to 4,096 bytes and segments BW500 writes into measured uncompressed
deliveries of up to 203 bytes. A wide level-4 delivery can carry two 89-byte
full frames and a closing frame of up to 86 bytes. Wide level-3 and level-2
full frames carry 47 and 22 bytes, respectively.
The optional CPAL bridge provides native capture and playback for isolated
audio routes. Station integration still needs qualified hardware playback
completion, PTT release, and the station admission and ownership contracts.

ARM64 preview installers and station guides are available for macOS,
Windows, Ubuntu, and Fedora. Windows, Ubuntu, and Fedora passed install,
launch, simulation, and removal checks. The extracted macOS package
executable ran successfully, and its bundled guide was checked. Native macOS
capture awaits an input endpoint.

\section{DSP implementation pipeline}
\label{sec:modem-dsp}

The standalone modem implementation reference expands the implemented
waveforms, trellis metrics, and bench recovery traces.
Its source and build instructions are in \path{docs/modem-implementation/}.

% Shared by the architecture white paper and the modem implementation reference.
\providecommand{\modemcite}[1]{}

\subsection{Audio ingress, continuity, and candidate detection}
The optional CPAL input backend requires native 48,000 samples/s and a selected
zero-based channel; it accepts float32, signed 16-bit, or unsigned 16-bit
formats and converts to normalized mono float samples. It selects one channel
rather than averaging potentially phase-opposed channels. No implicit
resampling is implemented. The callback accumulates 1,024-sample blocks and
uses a nonblocking bounded queue of 64 blocks. This represents approximately
1.365 s of mono input; it is a queue capacity, not a permitted DSP or keying
latency. Queue overflow, a malformed interleaved input block, or device failure
latches a terminal capture fault. The receiver must restart on a new continuous
stream instead of concealing dropped samples.

DSP runs outside the capture callback. The streaming receive worker validates
finite samples before advancing state, bounds processing steps to 480 samples,
and retains at most 30 s of mono history (1,440,000 float samples, about
5.76 MB of sample payload). Demodulation allocates additional bounded working
arrays, so the retained-history limit is not a whole-process memory limit.
Absolute half-open sample ranges survive chunk boundaries and ring-buffer
wrap. If a candidate began before retained history, it is marked truncated and
emits no partial frame. WAV reports count truncated energy candidates and
retain independently acquired pilot frames.

The energy detector accumulates 10 ms windows. A window is active when
\[
 \frac{1}{W}\sum_{n=0}^{W-1}x[n]^2 \geq 10^{-4},\qquad W=480.
\]
A 10 ms window holds 480 samples at 48 kHz. Square each normalized
sample so positive and negative excursions both count, then average.
A mean square of $10^{-4}$ corresponds to an RMS amplitude of 0.01,
so this gate opens at one percent of full scale.

A quiet window closes a candidate; end-of-stream flushes a partial window and
any open candidate. The 1\% full-scale RMS gate serves level-4 acquisition
with an external measurement profile and reports energy candidates for
diagnostics. The CLI scans level-1--3 pilots independently of this gate.
Energy candidates can include impulsive noise and unrelated signals.

\subsection{BW500 levels 1--3: tone-index geometry}
\modemcite{proakis}These levels encode data by selecting a frequency bin for each symbol. They
are not the level-4 differential-BPSK waveform played more slowly. The
implementation performs direct sine/cosine dot products at selected DFT bins;
it does not compute a full FFT for each symbol.

\begin{table}[htbp]
\centering\small
\begin{tabular}{@{}lrrrrrrr@{}}
\toprule
Level & $H$ & Bins & Columns & Refs. & $N$ & Coded & Payload\\
\midrule
1 & 2048 & 54--74 & 124 & 24 & 96 & 300 & 9 B\\
2 & 1024 & 27--37 & 228 & 24 & 200 & 612 & 22 B\\
3 & 1024 & 27--37 & 226 & 24 & 297 & 606 & 34 B\\
\bottomrule
\end{tabular}
\caption{Implemented index-mode geometry at 48 kHz. Columns exclude the initial lead; $N$ includes mode padding before whitening. Payload excludes the control byte and two-byte CRC.}
\label{tab:dsp-index-geometry}
\end{table}

The symbol rates are $48000/H$: 23.4375 symbols/s at level 1 and 46.875 at
levels 2/3. The synthetic transmitter adds two lead symbols at level 1 and four
at levels 2/3, all totaling 4,096 samples. Those bench leads use fixed in-band
tones; they do not implement the stock session's randomized lead stream.
There are three coded bits per data column. Subtracting the 24 reference
columns gives $100\times3=300$, $204\times3=612$, and $202\times3=606$
channel bits respectively.

Mode construction uses seeds 12,600,710; 7,910,701; and 14,955,604. One generator
stream produces the channel permutation, then the turbo permutation, followed
by bin/pilot draws. Level 1 replaces its generated base bins with the explicit
measured 124-entry array in the source. Levels 2/3 currently use generated
base allocations. All three therefore need qualification appropriate to their
actual construction, not a blanket claim that every allocation is independently
measured.

For each data column, the eight labels $(0,1,3,2,7,6,4,5)$ select offsets from
the column's base bin. Level 1 applies stride 2 modulo a span of 21; levels
2/3 apply stride 1 modulo 11. Fixed reference columns and their known bin
values allow synchronization without comparing decoded payload to an expected
message. The level-3 information block contains 37 complete bytes plus one
padding bit: 34 payload bytes, one control byte, two CRC bytes, and a zero
padding bit before whitening.

\landscapeplantuml{build/diagrams/Linkwise_DSP_Index}{Implemented BW500 level 1--3 receive pipeline. Pilot acquisition locates index frames in continuous audio. Soft demodulation and CRC validation follow timing and frequency selection.}{fig:dsp-index}

\subsection{Index-mode acquisition and soft demodulation}
\modemcite{proakis,meyr}The WAV and device receive commands use \path{modem::bw500::scan} to
search every 128 samples in a bounded rolling spectrum history. At each
epoch, the scanner scores the 24 reference columns against energy in
the occupied band and guard bins at offsets of 0 and $\pm$12 Hz. It
considers hypotheses with a mean pilot score of at least 0.65 and
refines timing in 16-sample steps. This lets a frame span brief drops in
energy or continuous receiver noise. The library also retains the
candidate-based decoder, which starts its search near $b+LH$ for burst
start $b$ and lead length $L$. The score threshold has no measured
false-alarm rate.

For a symbol of length $H$ and bin $k$, the unnormalized tone energy is
\[
 E_k=\left|\sum_{n=0}^{H-1}x[n]\exp(-j2\pi kn/H)\right|^2.
\]
The exponential supplies a reference that completes $k$ cycles over the
$H$ samples. Multiplying and summing aligns a matching tone so its
contributions add together. Taking the squared magnitude measures the
strength of that match independently of its starting phase. For example,
at 48 kHz with $H=1024$, bin 32 tests for a 1,500 Hz tone.

For timing refinement and soft decisions, the implementation divides by
$(H/2)\sum x[n]^2$, with an epsilon floor. For data, it measures the
eight mapped bin hypotheses. For each of the three bit positions, the soft
metric is the largest normalized power among labels with bit 1 minus the
largest among labels with bit 0. This max-based difference is a decoder input;
it is not a calibrated log likelihood, SNR, or BER estimate.
A positive value means the strongest tone compatible with bit 1 beat
the strongest tone compatible with bit 0. Values near zero give the
decoder little reason to prefer either bit.

Metrics are scattered through the inverse channel permutation into FEC order.
Only a CRC-valid result is returned. The rolling scanner suppresses overlap
with an emitted frame and continues to search for adjacent frames. Level 4
handles one frame per energy candidate. Continuous carrier and sample-clock
tracking remain open.

\subsection{Index-mode transmit synthesis}
The level-1/2/3 transmitter accepts a nonempty list of frames with exactly
9, 22, or 34 payload bytes per frame. It appends the control byte and byte CRC,
serializes MSB-first, pads to the mode's $N$, whitens, and turbo-encodes. For
each data column it gathers three channel-ordered coded bits into an integer
0--7 and selects the corresponding mapped tone. Reference columns emit their
fixed expected bins without consuming coded bits. A burst has one lead at its
start, followed by consecutive frame-column blocks; the lead is not repeated
between frames.

Each selected bin $k$ is rendered as $0.5\sin(2\pi kn/H)$ for $n=0,\ldots,H-1$.
The phase argument restarts for each symbol; integer-bin tones complete an
integer number of cycles over $H$ samples. The implementation concatenates
those symbol buffers directly. It does not add DATA-symbol pulse shaping,
resampling, a stock randomized acquisition stream, or a qualified transmitter
spectral mask. The fixed half-scale amplitude is a bench sample convention.
Reference-peer acceptance, occupied bandwidth, audio-route gain, and actual
RF power must be measured separately before treating this synthesizer as an
operational transmitter.

\subsection{Shared CRC, whitening, and interleaver generation}
\modemcite{koopman}Frame integrity uses CRC-16/GENIBUS: polynomial \texttt{0x1021}, initial register
\texttt{0xffff}, final complement, no reflection, and a big-endian two-byte
trailer. The check vector is \texttt{123456789} $\mapsto$ \texttt{0xd64e}.
The CRC covers the actual frame body, including its control field; mode padding
bits are outside the byte CRC. CRC checks establish frame integrity under this
code, not authenticated sender identity or recipient delivery.

Bits are serialized most-significant first, padded to the mode's information
size, and XORed with a freshly seeded whitening sequence. The 24-bit generator
obeys
\[
 s_{i+1}=(\texttt{0x43fd43fd}\,s_i+\texttt{0xc39ec3})\bmod 2^{24}.
\]
Each step multiplies the current state $s_i$ by the first constant,
adds the second constant, and keeps the lowest 24 bits. The remainder
operation is what makes the state wrap. The same starting state and
steps produce the same bit sequence at both ends of the link.

Whitening starts at $s_0=29318$ and takes the next state's high bit. A
generator bit of 1 flips the corresponding frame bit. A generator bit
of 0 leaves it alone. Applying the same mask again flips those bits
back. Permutations use
rejection sampling: draw an index, append only its first occurrence, and
continue until every index is present. The draw multiplies the normalized
state by the requested size with single-precision rounding, matching the
implementation's wire-oriented generator semantics. Changing numeric precision
or substituting a library shuffle can change the permutation.

The channel permutation and turbo permutation are different objects. If
$\pi_1$ maps transmitted position $j$ to a coded-bit position, TX gathers
$q_j=c_{\pi_1(j)}$ and RX scatters the corresponding soft metric back into
$c_{\pi_1(j)}$. The inner permutation $\pi_2$ changes the input order of the
second constituent encoder and must be inverted consistently when exchanging
extrinsic information. Interleavers belong to the mode/profile, not to the
agent's request parameters.
In the worked example, the channel permutation puts coded bit 33 in the
first transmitted position. The receiver puts the first soft decision
back into slot 33 before decoding. The turbo permutation acts earlier,
when the second encoder reads the information bits in a different order.

\subsection{Turbo encoder and iterative receive decoder}
\modemcite{berrou}The codec has two eight-state recursive systematic convolutional constituents.
Writing the three state bits as $(r_0,r_1,r_2)$ and the next information bit as
$u$, the implementation uses
\[
 a=u\oplus r_1\oplus r_2,\quad
 p=a\oplus r_0\oplus r_2,\quad
 (r_0',r_1',r_2')=(a,r_0,r_1).
\]
The symbol $\oplus$ means exclusive OR. First compute the new feedback
bit $a$ from the input and two stored bits. Then compute the parity bit
$p$ from $a$ and two stored bits. Finally, shift the memory so $a$
becomes the newest bit. Starting with an empty register and input 1,
the encoder outputs parity 1 and moves to state $(1,0,0)$.

Both encoders start at zero. The first consumes natural-order bits and the
second consumes turbo-permuted bits. Levels 1/2 transmit systematic and both
parity streams (rate 1/3); levels 3/4 transmit systematic bits and alternate
which constituent parity is retained (rate 1/2). Each constituent is terminated
with three steps using $u=r_1\oplus r_2$. The wire tail is twelve bits,
interleaving the two systematic/parity tail pairs at each of three steps.
Thus the implemented coded sizes are $3N+12$ or $2N+12$.
The factor of three or two counts the transmitted bits per information
bit before termination. For the worked level-4 frame, 368 information
bits give 736 systematic and parity bits, then the tail brings the
total to 748.

\landscapeplantuml{build/diagrams/Linkwise_DSP_FEC}{Implemented shared whitening, constituent coding, tail and permutation structure. Receive decoding uses a hard-decision CRC fast path followed by at most eight max-log-MAP iterations.}{fig:dsp-fec}

\modemcite{bcjr,robertson}The receiver first makes hard decisions on systematic metrics and checks the
frame CRC. If that fails, it constructs the two terminated trellises, treating
punctured parity as a zero-information metric. A forward/backward
\emph{max-log-MAP} pass computes extrinsic information for each constituent.
The implementation uses maxima rather than log-sum-exp correction; it must
not be described as an exact log-MAP decoder. Positive input metrics favor
bit 1. The two constituents exchange interleaved/de-interleaved extrinsic
information for at most eight iterations, checking de-whitened frame integrity
after each iteration and stopping on a valid CRC.

This implementation does not maintain cross-retransmission soft combining.
Its metrics have heuristic scales and its tests do not establish coding gain,
waterfall curves, or residual undetected-error probability under HF fading.
Those require independently generated vectors and controlled channel sweeps,
including offset, burst noise, clipping, and long erasures.

\subsection{Level-4 frame geometry and external measurement profile}
The level-4 byte container is 46 bytes: 43 payload bytes, one control byte,
and two CRC bytes. Its $N=368$ information bits produce 748 coded bits. The
channel and turbo permutations start from seed 6,437,280. Pilot generation
continues that stream after the two permutations, choosing 18 distinct columns
per band with $5+\lfloor R\,386\rfloor$.

The implemented waveform stores 393 columns per band: one reference followed
by 392 body columns. Removing 18 pilot positions in each band yields
$2(392-18)=748$ differential data cells. This is the retained single-frame
implementation geometry, not a claim that every stock burst has exactly that
many emission slots. Its reference centers start after nine 512-sample slots.
Band centers are 1,350 and 1,650 Hz, the symbol hop is $H=512$ samples
(93.75 symbols/s), and the upper band is staggered by 256 samples. Unlike the
index ladder, these carriers are not selected integer bins of a symbol FFT.

The phase reference consists of 786 measured points on a 128-phase lattice,
393 per band. Each point is $G_{b,t}=\exp(j2\pi k_{b,t}/128)$. A shared measured
513-sample real pulse covers offsets $-256$ through $+256$. The profile loader
requires all reference positions, rejects duplicates and out-of-range phases,
checks pulse length and finite/bounded taps, and limits JSON size to 128 KiB.
Those structural checks do not prove that an arbitrary supplied profile matches
a real modem. Reproducible tests use SHA-256-pinned upstream measurements,
extracted without executing downloaded source. Cached measurement sources
retain their own attribution/license metadata and are not embedded in Rust.

The phase index $k_{b,t}$ selects one of 128 equally spaced directions
around a circle. For example, index 32 represents a quarter turn. The
transmitter and receiver use opposite rotations so they can compare
successive DATA symbols after accounting for the known phase pattern.

\landscapediagram{\begin{tikzpicture}[x=1.02cm,y=1cm,font=\small,
  axis/.style={ovgray,-{Stealth[length=2mm]}},
  data/.style={ovblue,fill=ovblue},
  pilot/.style={ovgold!60!black,fill=ovgold},
  reference/.style={ovteal!70!black,fill=ovteal}]
  \node[anchor=west,font=\bfseries] at (-2.7,6.15) {BW500 level-4 symbol centers};
  \node[anchor=west] at (-2.7,5.55) {$n_{b,j}=4608+512j+256b$, with band index $b\in\{0,1\}$ and column $0\le j\le392$};
  \foreach \x in {0,1,...,17}
    \draw[ovgray!20] (\x,1.05) -- (\x,4.95);

  \draw[axis] (0,3.9) -- (17.65,3.9);
  \draw[axis] (0,2.45) -- (17.65,2.45);
  \node[align=right,anchor=east] at (-0.35,3.9) {Lower band\\1,350 Hz};
  \node[align=right,anchor=east] at (-0.35,2.45) {Upper band\\1,650 Hz};
  \draw[reference] (-0.09,3.81) rectangle ++(0.18,0.18);
  \draw[reference] (0.91,2.36) rectangle ++(0.18,0.18);
  \foreach \j in {1,...,8} {
    \fill[data] ({2*\j},3.9) circle[radius=2.4pt];
    \ifnum\j=7
      \draw[pilot] (15,2.59) -- (15.14,2.45) -- (15,2.31) -- (14.86,2.45) -- cycle;
    \else
      \fill[data] ({2*\j+1},2.45) circle[radius=2.4pt];
    \fi
  }
  \node[above,font=\scriptsize] at (0,4.11) {$j=0$};
  \node[above,font=\scriptsize] at (2,4.11) {$j=1$};
  \node[above,font=\scriptsize] at (15,2.59) {$j=7$, first upper pilot};
  \draw[ovblue,{Stealth[length=1.4mm]}-{Stealth[length=1.4mm]}]
    (0,4.7) -- (2,4.7) node[midway,above,font=\scriptsize] {512-sample hop};
  \draw[ovteal!70!black,{Stealth[length=1.4mm]}-{Stealth[length=1.4mm]}]
    (0,3.15) -- (1,3.15) node[midway,above,font=\scriptsize] {256};

  \draw[axis] (0,1.05) -- (17.65,1.05);
  \foreach \x/\sample in {0/0,2/512,4/1024,6/1536,8/2048,10/2560,12/3072,14/3584,16/4096}
    \draw[ovgray] (\x,1.05) -- ++(0,-0.11) node[below,font=\scriptsize] {\sample};
  \node[anchor=west] at (0,0.34) {Samples relative to the lower-band reference, $n-4608$};

  \draw[reference] (0,-0.55) rectangle ++(0.18,0.18);
  \node[anchor=west] at (0.32,-0.46) {Reference};
  \fill[data] (3.05,-0.46) circle[radius=2.4pt];
  \node[anchor=west] at (3.27,-0.46) {Data};
  \draw[pilot] (5.1,-0.32) -- (5.24,-0.46) -- (5.1,-0.60) -- (4.96,-0.46) -- cycle;
  \node[anchor=west] at (5.34,-0.46) {Pilot};
  \node[anchor=west,font=\scriptsize] at (7.0,-0.46)
    {18 pilots per band. The first lower-band pilot is at $j=24$.};
\end{tikzpicture}
}{BW500 level-4 symbol centers on the 48 kHz sample axis. The plotted excerpt uses the implemented pilot allocation. Each marker locates the center of a 513-tap pulse extending 256 samples on either side.}{fig:dsp-level4-symbol-timing}

\subsection{Level-4 transmit synthesis}
After whitening, turbo encoding, and channel gather, coded bits fill band 0's
data positions in time order, then band 1's. With initial polarity
$v_{b,0}=+1$, a coded 1 negates polarity, a coded 0 preserves it, and a pilot
repeats the preceding polarity. Each complex emitted cell is
$c_{b,t}=v_{b,t}\overline{G_{b,t}}$. The real audio is a sum of pulse-shaped
carriers:
\[
 x[n]=\frac14\sum_{b=0}^{1}\sum_t p[n-n_{b,t}]
       \Re\{c_{b,t}\exp(j2\pi f_b n/48000)\},
\]
where $n_{b,t}=9H+tH+bH/2$, $f_b\in\{1350,1650\}$, and $p$ has 513 taps.
To compute one output sample $x[n]$, center each symbol's pulse on
$n_{b,t}$, evaluate the carrier at sample $n$, and add the real
contributions from both bands. The pulse determines how much each nearby
symbol contributes. The center formula allows nine hops before the
first reference, advances one hop per column, and offsets the upper
band by half a hop.

The factor $1/4$ is the current bench amplitude scale, not a calibrated radio
power command. A waveform's audio amplitude does not replace CAT power limits,
route qualification, or PTT authority.

The synthetic buffer contains $9H+393H+1=205825$ samples, approximately
4.288 s at 48 kHz. The initial slots do not contain a generated randomized
acquisition lead. This output length is a local implementation fact and must
not be substituted for a measured reference modem's PTT interval, framing
duration, or turnaround requirement. Actual device playback and RF keying are
not performed by the file-only synthesis tool.

\subsection{Level-4 timing lock, demodulation, and validation}
\modemcite{proakis,meyr}The receiver mixes the real input down at each nominal carrier using
$x[n]\exp(-j2\pi f_b n/48000)$. At a trial symbol center it multiplies the
513 surrounding complex samples by the measured real pulse and sums them.
It then multiplies by $G_{b,t}$ to remove the known phase reference. If the
result is $z_{b,t}$, the differential statistic is
\[
 d_{b,t}=\Re\{z_{b,t}\overline{z_{b,t-1}}\}.
\]
This dot product compares consecutive symbols after removing the known
reference phase. It tends to be positive when their polarities agree
and negative when they are opposite. A DATA bit 1 flips polarity,
so the receiver negates this statistic to give a positive score for 1.

A correct pilot transition should preserve polarity. Timing search maximizes
the average normalized differential dot product over 36 pilot positions,
first in 32-sample steps and then in two-sample steps around the best epoch.
A score below 0.65 is rejected. The search is bounded by the candidate, frame
geometry, available samples, and one-hop end tolerance; it does not scan an
unbounded recording for a matching payload.

Data positions produce $-d_{b,t}$ so positive metrics favor a transmitted 1.
The 748 values are scaled by six divided by their mean absolute magnitude,
then passed through channel scatter and the shared turbo/de-whitening/CRC
chain. This normalization is not an estimate of a physical noise variance.
No valid CRC means no frame. A valid frame is interpreted in protocol context:
a setup frame can identify a caller; a DATA container can carry payload and
control; neither is an unconditional link-up event.

\landscapeplantuml{build/diagrams/Linkwise_DSP_Level4}{Implemented level-4 transmit and receive transforms. Measured reference phases and pulse taps are external profile inputs; the single-frame synthetic round trip does not qualify reference interoperability.}{fig:dsp-level4}

The level-4 decoder currently allocates complex mixed buffers and performs
repeated direct matched-pulse sums. It is a bench implementation, not a
profiled real-time DSP scheduler. An optimized runtime may cache oscillators,
reference terms, pulse weights, and mode tables, and bound search work per
candidate, but must preserve sample epochs, metric sign, permutation order,
and CRC outcomes against golden fixtures before replacing these kernels.


\subsection{Proposed production DSP extensions and work budgets}
The following is a proposed engineering path, not behavior implemented by the
current decoder. A production receiver needs explicit acquisition/tracking
states, measured latency budgets, and loss-of-lock rules rather than extending
the bench thresholds by assumption. Mode tables, pulse taps, and reference
terms should be validated and cached before a session; source loading,
allocation-heavy analysis, persistence, and model calls stay outside the
capture and playback deadlines.

\begin{enumerate}[leftmargin=1.6em]
  \item \textbf{Front-end conditioning and health.} Measure route noise,
        clipping, DC offset, stream discontinuities, and usable passband.
        Select channel filters and gain normalization from qualified profiles.
        An adaptive candidate detector needs a recorded noise estimator,
        hysteresis, duration limits, and a false-alarm/miss test population.
        Gain changes must not erase sample epochs or turn audio energy into
        protocol evidence.
  \item \textbf{Coarse acquisition.} Search a bounded frequency/timing bank
        against references before attempting expensive FEC. Record competing
        mode scores and ambiguity rather than treating the first energetic
        interval as a known modem. Reuse selected-bin correlations or an FFT
        bank only after latency measurements establish the benefit and fixture
        tests preserve the selected bins and normalized scores.
  \item \textbf{Carrier tracking.} After correct reference de-rotation and
        under a locally constant channel, repeated pilot transitions provide
        a candidate residual-offset estimator:
        \[
          \widehat{\Delta f}=\frac{48000}{2\pi H}
          \arg\!\left(\sum_{(b,t)\in\mathcal P}
          w_{b,t}z_{b,t}\overline{z_{b,t-1}}\right).
        \]
        Each pilot pair measures the phase change over one symbol hop.
        The set $\mathcal P$ identifies the pilot positions and the
        weights $w_{b,t}$ determine how much each pair contributes.
        The angle of the combined result gives radians per hop.
        Multiplying by $48000/(2\pi H)$ converts that angle to hertz.
        For example, a quarter-turn of phase per 512-sample hop
        corresponds to about 23.44 Hz.

        This is a proposed pilot estimator, not a performance result. At
        $H=512$ its phase observation repeats every 93.75 Hz, requiring coarse
        branch selection and outlier/low-energy rejection. Fading and channel
        variation must be included in qualification before closing a loop.
  \item \textbf{Sample-clock and channel tracking.} Track fractional symbol
        epochs and their drift over long bursts; qualify interpolation or
        resampling with known offsets rather than slipping whole samples
        silently. Assess pulse distortion and multipath-induced intersymbol
        or interband interference. Any adaptive equalizer needs training,
        convergence, divergence, and reset criteria against recorded HF
        channel cases. A capture gap invalidates continuity and tracking state.
  \item \textbf{Metric and deadline qualification.} Compare heuristic
        power/differential metrics with noise-aware likelihood models; validate
        sign, clipping, erasures, and decoder response before changing scales.
        Bound candidate search, queue occupancy, FEC iteration count, and
        cancellation latency. Measure worst-case and high-percentile service
        times on each target under simultaneous capture, decode, client I/O,
        and audit load. Drop optional diagnostics before compromising sample
        continuity; if continuity or a safety deadline cannot be maintained,
        terminate the affected stream/session explicitly.
\end{enumerate}

The current turbo decoder performs forward/backward passes on two eight-state
trellises, so work scales linearly with information size and iteration count;
forward metric storage is also linear in block size. This algorithmic bound
does not establish a real-time scheduling bound. Index acquisition repeatedly
correlates reference windows, and level-4 trial epochs repeatedly evaluate
matched pulses; both can dominate a failed candidate. Kernel benchmarks must
therefore include failed acquisition and failed CRC/FEC paths, not only the
fast hard-decision success case.

A DSP qualification record should identify source audio/hash, route and sample
format, mode/profile hash, impairment model and seed, expected frame/payload,
acquisition epoch, frequency/clock estimates, decoder iterations, terminal
CRC result, queue high-water marks, and elapsed work. BER requires known
reference bits; a CRC failure alone does not measure BER. Fixture-defined
frame-error rates, energy-candidate rejection rates, and link/session outcomes
need separate denominators. Actual output presentation latency and audio/PTT
lead/tail behavior require device-loopback measurements in addition to these
software records.

\subsection{Control bursts, caller setup, and handshake evidence}
Control bursts use a separate 2,048-sample tone grid at 48 kHz. Request and
response tones are deterministic functions of the known called station;
their payload uses the BW500 even-bin alphabet 50--76 while fixed preambles
may use other bins. The request contains ten preamble plus 31 payload symbols;
the response contains eight plus 15; the level-4 confirm contains one plus 15.
Connected ACKs contain eleven two-tone symbols, with a fixed four-symbol
preamble and a caller-keyed tail. The control renderer uses a 32-sample overlap
and cosine edge shaping. These control bursts do not pass through the DATA
turbo/interleaver chain.

Recognizers compare the complete expected keyed sequence near the estimated
burst start. Single-tone matching requires an average normalized tone score
of at least 0.8 and a weakest-symbol score of at least 0.30; connected-ACK
matching requires total pair score of at least 0.8 per symbol and at least
0.15 for each individual expected tone. These are bench thresholds, not
field-qualified receiver sensitivities. The complete ACK tail is checked;
shared preamble energy alone cannot establish acceptance. Request/response
matching cannot discover an unknown caller because those bursts are keyed to
the called station. Caller identity arrives in the CRC-protected setup frame.

The setup codec packs up to six base callsign characters into six-bit symbols,
retains an SSID in byte 5, records capability/request fields, includes the high
byte of the full-callsign CRC, and finishes with a frame CRC. Current native
audio establishment is restricted to level 4. Its state machine expects request,
response, setup, connected ACK, then confirm; wrong-phase or wrong-peer bursts
cannot advance it. Unique local transmit tokens reject stale completion
callbacks. Establishment is emitted only after the relevant confirm is
received or actually transmitted, rather than when a waveform is merely queued.

The current two-role integration test supplies transmit completion after the
synthetic audio has been exchanged. A production device owner must provide
that event only after actual presentation and PTT release. Software buffer
consumption is not proof that the final sample reached the output device.
Clocked playback, transmit/receive turnaround, and qualified PTT readback
remain required integration work.

\subsection{DATA framing, ARQ ownership, and ambiguous answers}
The DATA container codec is separate from waveform decoding. At level 4 the
body has 44 bytes, at level 3 it has 35, and the final body byte is a control
field. A full frame uses the preceding bytes for payload. The supported short
form places payload, a \texttt{0x14} marker, the caller CRC high byte, zero
fill, a next-receive-level byte, and \texttt{0x82} before the CRC. Unsupported
or ambiguous forms are rejected. The implemented short builder reserves four
body bytes for this trailer; additional published short/compressed forms are
not silently treated as supported.

The BW500 decoder also accepts the measured level-4 \texttt{0x02}
control form. In setup, it validates the CRC, caller identity and
requested level. In DATA, it requires the caller-keyed trailer and valid
padding and rate fields. Stock Winlink Express used this form to send
the B2F end-of-proposals bytes after storing the received message.

Tests reconstruct the documented 200-byte delivery from four level-4 payloads
of 43 bytes and one closing level-3 payload of 28 bytes, using the documented
full fields \texttt{95}, \texttt{91}, \texttt{8d}, and \texttt{81}. This verifies
the container interpretation for that example, not a general sender sequencing
law. The full-frame field remains a supplied parameter because its evolving
wire semantics need independent qualification. Appendix~\ref{app:bw500-wire}
works through these bytes, their checksums, and the session exchange.
Compression TEXT/FILES,
delivery-level repacketization, and negotiated rate changes are not completed
by the container codec.

The event-level ARQ sender owns bounded queued blocks and exactly one
outstanding block. DATA transmit completion opens a receive deadline. A
validated delivered reply retires that block once; a validated not-delivered
reply permits a retry within budget. A missing reply is ambiguous: the peer
may have accepted DATA while its answer was lost. The timeout therefore emits
a state-query action rather than blindly repeating DATA. Query transmission
has its own completion token and watchdog. Query/retry exhaustion, transmit
failure, or missing/late playback completion faults the sender and preserves
unresolved bytes. Explicit abort discards the queue.

These local tokens bind callbacks to an outstanding transmission/window; they
are not transmitted sequence numbers or a solution to on-air duplicate
suppression. The station/native session owner must connect qualified reply
recognition, query encoding, turn ownership, negotiated rate, and byte
accounting. The sender's queue currently holds supplied blocks; its byte count
must not be exposed as Winlink delivery progress or raw-host \texttt{BUFFER}
semantics without that integration.

The virtual host issues keyed state queries after missing BW500
intermediate and closing acknowledgments. A positive answer retires the
acknowledged bytes. A negative answer can trigger a level-3 retry of the
outstanding DATA, preserving any acknowledged prefix. With crossed turn
requests, Linkwise can grant a repeated stock caller request and resume
queued DATA after the caller hands back the turn. A short continuation
can retain level 3 after a level-3 close. A host \texttt{DISCONNECT}
received while DATA is queued waits for the queue to drain before
sending the disconnect waveform.

For the wide modes, the virtual host tracks the outstanding waveform and
the byte offset at the active rate. It can retry the unacknowledged suffix
at level 3 after a keyed rejection, use level 2 after a missing-DATA query
response, and return to level 4 when the peer requests the supported
increase. These transitions cover measured frame layouts. General rate
negotiation and repeatable recovery across the qualification matrix
remain open.

\landscapeplantuml{build/diagrams/Linkwise_Modem_Recovery}{Implemented event-level ARQ sender behavior and remaining wire integration. Missing answers request state; only a validated negative reply permits DATA retransmission.}{fig:modem-recovery}

\subsection{Execution isolation, keying, and failure containment}
The proposed runtime separates four execution responsibilities. The input
callback only converts and queues audio; DSP workers detect, synchronize,
decode, and synthesize outside that callback; one native session actor owns
protocol transitions, outstanding bytes, and monotonic deadlines; the station
execution service owns the radio lease and client/modem operation. The PTT
arbiter and watchdog retain independent release authority. MCP queries read
snapshots and append commands to bounded application queues; slow model
responses, JSON serialization, audit storage, and provider fetches never run
inside an audio callback or protocol deadline.

A burst request must identify the session, active lease, configuration
revision, permitted route, and bounded duration. The arbiter validates those
against policy and interlocks before key-up, arms release limits, and uses a
qualified driver. Playback completion and release observation are separate
facts. Timeout, device loss, capture discontinuity, lease loss, abort, or an
unconfirmable release cannot be converted to a successful session. A release
request still proceeds when audit persistence fails; the fault is latched and
reported separately. The physical emergency stop remains independent of
software execution and of either AI host.

\landscapeplantuml{build/diagrams/Linkwise_Modem_Keying}{Proposed local playback, keying, completion, and watchdog sequence below an asynchronous agent request. This integration is not yet implemented; no MCP round trip is allowed inside a symbol or ACK deadline.}{fig:modem-keying}

\subsection{Verification matrix and implementation exit gates}
The 9 November run of \path{tools/test-modem.py --audio-device} passed 41 tests:
37 library tests, one replay integration test covering five independently
published level-1/2 excerpts, and three synthetic level-4/audio integration
tests. The reference clips include a two-frame burst and retain upstream
provenance and SHA-256 checks. The level-4 tests cover setup and DATA round
trips, local two-role establishment, chunked streaming, and corruption
rejection. Formatting and Clippy checks also passed. These are dated evidence,
not a promise that subsequent source changes have identical results.

On 10 November the fixture-enabled suite passed 68 tests on macOS
with the audio feature enabled. The pilot acquisition regression
covered consecutive frames in continuous interference, chunk independence,
and duplicate suppression. Formatting and Clippy passed. The offline
receive checks used the full demodulated recordings and supplied no burst
boundaries. Capture settings, hashes, and reproduction commands are in the
10 November entry of \path{lab-notes/lab-notes.tex}.

The protocol notes and reference measurements are pinned to upstream revision
\path{952f79245240132f4af209afbad1f4b2ba256613}; implementation details should
be checked against the named Rust modules and \path{rust/docs/modem.md}.
The notes, fixtures, and measured sources have distinct attribution/license
metadata; the cached measurement inputs are not proprietary vendor assets
embedded in this repository. The public protocol/reference source is listed
in Reference~\ref{ref:modem-protocol}. Source-level traceability does not
replace empirical qualification.

\paragraph{Deterministic waveform regression and application delivery.}
The 12 November 2026 BW500 suite passes packed PRBS9, PRBS15, alternating
bits, and all-zero and all-one payloads through Linkwise's real-valued
48~kHz waveform path. Levels 1--3 recover the exact DATA frames through
attenuation and a fixed 1~ms echo. A separate level-4 test reassembles
the 64-byte PRBS9 fixture from full and closing DATA frames using the
measured waveform profile. The virtual daemon bench also recovers
PRBS9 payloads at BW500, BW2300, and BW2750, with controlled DATA
corruption or acknowledgment loss in selected cases.

On the mixed BW500 bench, both Winlink Express peers sent and received
an email with a PRBS9 attachment in each of two sessions. Linkwise and
stock VARA HF 4.9.1 each initiated one session. Both peers disconnected
cleanly and the four received attachments matched their sources byte for
byte and by SHA-256. Separate stock-audio replays cover the initial caller
poll and rate-down recovery in both caller roles. These measurements
cover frame decoding and delivery after ARQ. Measuring raw bit errors
requires aligned receiver decisions, and assessing fading tolerance
requires a calibrated channel model. Fixture construction, commands, and
run scope appear in the 12 November entry of \path{lab-notes/lab-notes.tex}.

\begin{longtable}{@{}p{0.24\textwidth}p{0.70\textwidth}@{}}
\caption{Modem evidence and remaining qualification gates.}\label{tab:modem-evidence-gates}\\
\toprule
\textbf{Gate} & \textbf{Evidence and remaining checks}\\
\midrule
\endhead
Codec correctness & Independent vectors for whitening, permutations, FEC tails,
metric sign, frame fields and mode geometry; expand beyond self-generated
round trips and the level-1/2 corpus.\\
Off-air receive & One CRC-valid BW500 level-2 frame from the full Lake
Merwin KiwiSDR recording on 7.105 MHz USB. The published off-air exchange
yielded three CRC-valid frames (levels 2, 3, 2). Both replays reproduced
previously verified wire-body hashes.\\
Acquisition controls & Two stock-derived frames in recorded Kiwi noise
decoded automatically at +6 dB. The Maui noise-only replay emitted zero
frames. These checks cover the tested recordings.\\
Receiver robustness & Measure live frame recovery on known transmissions.
Characterize offset, fading, interference, clipping, sample-clock error,
false alarms, and processing latency.\\
Reference acceptance & With stock VARA HF 4.9.1, the 11 November frozen
build \texttt{6a4dc89e} passed one BW2300 and two BW2750 sender-recovery
runs with 412 exact bytes, two-ended disconnect and a same-socket
reconnect. Other runs stalled during setup or retry. The three-pass gate
remains open.\\
Session reliability & The frozen build passed two final-ACK-loss runs and
one intermediate-ACK-loss run at BW500, with 406 exact bytes and two-ended
disconnect in each pass. Setup and DATA timeouts interrupted the sequences.
The earlier three-pass results apply to their recorded builds.\\
Station safety & Qualify actual device presentation, turnaround, exclusive
lease, permit expiry, duty limits, watchdog release and independent stop on
an inhibited or conducted setup before operational RF use.\\
Host and Winlink & Mixed BW500 Pat runs passed in both directions with a
matching attachment SHA-256, exit status zero and two-ended disconnect.
Mixed Winlink Express 1.8.6.0 and Pat 1.0.1 peer-to-peer transfers stored
attachments with matching SHA-256 in both recipient mailboxes. The
Pat-to-Express sender-side close and gateway/CMS qualification remain open.
Two Winlink Express BW500 sessions with Linkwise and stock VARA verified
two-way PRBS9 attachment delivery and clean disconnect with each modem
as caller.\\
\bottomrule
\end{longtable}

The earlier i386 Debian/Wine bench on the ARM64 Mac stopped after an
illegal-instruction exception during COM registration and a WineDbg
attach error. The later Windows UTM bench provided a working stock peer.
The November 11 lab chronicle records the stock Windows UTM runs, the
Express message export, and installer checks. Its evidence manifest links
the measured results to the captured audio and command transcripts.
The repeat counts, build hashes, and failure captures are indexed in
\path{docs/testing/vara-interoperability.md}.
Physical transmit control and operational RF qualification remain open.


\section{Deployment profiles}

\begin{description}[leftmargin=1.45in,style=nextline]
  \item[Shadow qualification] The agent observes and makes simulated/shadow decisions; no transmitter actuation is enabled. This is a pre-deployment qualification stage, not an operator-selected runtime mode.
  \item[Autonomous operation] After commissioning and station-specific safety qualification, the agent has full locally enabled CAT/PTT authority and manages sessions, retries, and recovery automatically within the signed profile. Human runtime actuation is limited to emergency stop.
\end{description}

The autonomous profile is not activated until all readiness checks pass. During operation, policy faults cause automatic recovery or a latched safe state; only the independent human emergency stop is exposed as an operator actuator. Model/provider selection must not alter station-level limits.

The requirements and verification strategy set the gates for implementation.
The roadmap orders work through observation, policy-controlled actuation,
simulated/shadow qualification, and finally autonomous operation after station
safety acceptance.

\section{Implementation roadmap}

\begin{enumerate}
  \item \textbf{Foundation:} identify station profile, supported radio, and documented client interfaces; preserve evidence ledger.
  \item \textbf{Observability:} read-only Hamlib, modem, and gateway adapters; typed state and local audit records.
  \item \textbf{Inventory and history:} documented channel-source adapter, distinct gateway/channel counts, idempotent session/transfer events, and explicit metric coverage.
  \item \textbf{Candidate analysis:} source-normalized evidence, ranking/explanations, simulated test fixtures.
  \item \textbf{Preparation:} exclusive lease, configuration writes, readback, and mismatch handling; no transmission.
  \item \textbf{Session supervision:} supported Winlink client adapter, separate authorization, watchdog, timeouts, abort and recovery.
  \item \textbf{Modem interoperability:} build waveform analyzer and receivers/transmitter from verified captures; test against unmodified VARA in controlled conditions.
  \item \textbf{Adaptive operation:} learn from historical outcomes and enable only bounded automation that has passed station-specific review.
\end{enumerate}

\section{Risks and open decisions}

\clearpage
\begin{landscape}
\thispagestyle{plain}
\captionof{table}{Principal architecture risks and proposed dispositions.}
\label{tab:risks}
\begin{tabularx}{\linewidth}{@{}p{2.0in}X p{2.0in}@{}}
\toprule
\textbf{Risk/question} & \textbf{Impact} & \textbf{Disposition}\\
\midrule
Winlink client automation surface & Determines whether session orchestration can be robustly automated & Select/document API per client \\
Radio capability differences & PTT, power, bandwidth, and signal fields vary by hardware/backend & Capability-driven adapter; qualify models \\
Propagation uncertainty & Predictions and spots do not prove VARA reachability & Calibrate with session history and expose uncertainty \\
One-radio capacity & Cannot maintain concurrent frequencies/links & Optimize next session; multi-radio later \\
Unknown VARA behavior & Protocol implementation risks false assumptions & Capture-led clean-room research \\
Local legal/station limits & Incorrect configuration can create harmful transmissions & Operator-owned profile and hardware protections \\
External data licensing/availability & May constrain automation, retention, or redistribution & Review per source; graceful degradation \\
\bottomrule
\end{tabularx}
\end{landscape}
\clearpage

\section{Conclusion}

Linkwise's contribution is not a claim that AI can see through an uncertain
ionosphere. It is a disciplined way to connect existing radio-control,
modem, Winlink, and propagation software; assemble gateway, station, and
session evidence; compare available routes; and learn from measured outcomes.
The agent autonomously manages the configured station through mediated tools,
while deterministic local policy and watchdogs enforce the preinstalled
authority. During runtime the human has read-only telemetry and an independent
emergency stop. This architecture offers a path to autonomous communications
only after each adapter, control path, and operating limit has been validated.

\appendix
\section{Extended propagation model evidence}\label{app:model-evidence}
This appendix records model-comparison evidence that informs the
calibration limits discussed in the main text. Findings are specific
to their datasets, models, and evaluation methods.

Published comparisons demonstrate why the service must expose model scope and
uncertainty. In Walden's 23-month study of nine 5.29 MHz midlatitude NVIS
beacon circuits, prediction residuals varied with season and path. The
reported overall RMS difference was 7--15 dB for VOACAP and 9--16 dB for
ASAPS, while the study also found signals during some periods when the
models' support probabilities were zero [37]. In a 2026 global model
comparison using the selected CCIR measurement data and a common evaluation
procedure, REC533 had the lowest overall dispersion among the three models;
ICEPAC and VOACAP showed larger, condition-dependent errors, especially in
some long-distance, high-frequency, and low-latitude cases [39]. These are
results for the papers' datasets and methods, not a ranking that can be
transferred directly to Linkwise's RMS/VARA operating population. A forecast
should be calibrated and stratified by path class, season, solar/geomagnetic
regime, and operating configuration before it changes live candidate order.

Short-term forecasting of an ionospheric state is a distinct research problem
from predicting a particular oblique HF circuit. Moskaleva and Zaalov describe
a simplified model that forecasts F2-layer critical frequency and electron
density profiles one day ahead and compares the forecasts with ionograms [38].
Such profile skill can inform path models, but does not by itself establish
forecast accuracy for received field strength, channel availability, or
successful application transfer. Validation should be done at each intended
target metric using observations that represent that target; recent work
demonstrates evaluation of ionospheric models against technologically
relevant metrics including HF links [44].


\section{Template survey and document lineage}
The specified \urlref{https://github.com/8r4n/acq-latex-templates}{acq-latex-templates repository} was cloned and surveyed at commit \texttt{fef160e}. The Software Architecture Description (SAD) template draws on ISO/IEC/IEEE 42010: stakeholders and concerns, viewpoints and views, correspondences, and decisions with rationale. The Software Design Description (SDD) and Software Requirements Specification (SRS) templates add component-level design, execution behavior, interface and data design, traceability, and qualification methods. Those structures informed this paper and its companion design documents.

The acquisition-specific \texttt{dodacqdoc} class and marking conventions were not adopted: Linkwise is not a DoW acquisition deliverable. This paper uses a self-contained \texttt{article} layout. The assessment below is deliberately candid: a readable system whitepaper is not, by itself, a complete SAD plus SDD package.

For readability, every figure and table in this paper is presented on its own
landscape page with a descriptive caption. Long tables may continue onto
additional landscape pages, repeating their column headings.

\section{Architecture coverage assessment}\label{app:template-alignment}
The paper covers the major system-level questions an architect should answer,
but it does not claim full SAD/SDD conformance. The templates provide a
coverage checklist, not a reason to force every detailed design artifact into
one narrative. Companion documents hold the deeper views and requirements;
The companion SDD specifies layer responsibilities, component contracts,
principal execution flows, state/concurrency rules, persistence semantics,
and deployment evolution. Formal view correspondences, final
interface-control definitions, physical schema, implementation timing
budgets, and as-built qualification evidence remain future work.

\clearpage
\begin{landscape}
\thispagestyle{plain}
\begin{longtable}{@{}p{1.45in}p{1.25in}p{4.85in}@{}}
\caption{Coverage of software architecture and design elements from the surveyed SAD, SDD, and SRS templates.}\label{tab:template-coverage}\\
\toprule
\textbf{Template element} & \textbf{Assessment} & \textbf{Coverage or remaining work}\\
\midrule
\endhead
Purpose, scope, system overview & Covered & The abstract, executive summary, purpose, goals, and non-goals define the system and its proposed status. \\
Stakeholders and concerns & Covered at system level & The stakeholder table links primary concerns to operational, functional, deployment, safety, data, and verification views. Detailed stakeholder analysis belongs in the SAD. \\
Architecture viewpoints and views & Partial & The paper names principal viewpoints and includes context, operating-loop, component, deployment, evidence, and safety material. Formal model conventions and a complete view catalog belong in a fuller SAD. \\
Consistency and correspondences & Partial & Key allocations are stated here and mapped in the SAD cross-view table. A formal correspondence register with explicit rules and dispositions for inconsistencies remains needed for strict SAD completeness. \\
Decisions, alternatives, rationale & Covered at summary level & This paper compares principal alternatives and consequences. The SAD maintains the decision log; each record should include concern IDs, status history, and alternatives. \\
Components and execution concept & Covered at design-summary level & Layer ownership, component contracts, command sequencing, state transitions, concurrency, and recovery are described here and specified further in the companion SDD. Thread priorities, timing budgets, and per-component unit algorithms remain implementation work. \\
Interfaces and data design & Substantially covered; finalization pending & Typed port semantics, provenance-bearing observations, authorization handles, and event ownership are described in the paper and SDD. A finalized interface-control baseline and physical database schema are not yet present. \\
Requirements traceability and qualification & Partial & Representative requirements and methods appear here. The requirements document contains the fuller use-case, requirement-to-component, and qualification baseline; final bidirectional verification traceability remains a release task. \\
Security, safety, deployment, risks & Covered at proposed-system level & The paper identifies local control, PTT gates, trust concerns, deployment profiles, and major risks. Detailed hazards, controls, and failure handling are in the companion safety/security/reliability document. \\
As-built design and configuration record & Not applicable yet & This is a proposed architecture, not an as-built design. Implementation baselines, validated interfaces, timing/resource budgets, and final TBD values must be recorded before qualification. \\
\bottomrule
\end{longtable}
\end{landscape}
\clearpage

Overall, this whitepaper and its companion documents form a proposed
architecture-and-design baseline, not an as-built implementation record or a
claim of complete template conformance. The highest-value follow-up is to
formalize SAD view correspondences and decision records, stabilize and version
the interface contracts, and qualify the safety supervisor and adapters
against simulation and controlled hardware tests. Provisional choices remain
explicit so they are not mistaken for implementation facts.

\section{References}
\begin{enumerate}[label={[\arabic*]}]
  \item Hamlib, “Network Device Control,” \url{https://github.com/Hamlib/Hamlib/wiki/Network-Device-Control}.
  \item VOACAP Online, “User Manual,” \url{https://www.voacap.com/2023/documents/VOACAP_Online_User_Manual_16_July_2024.pdf}.
  \item Winlink, “Live System Information / RMS Channels,” \url{https://winlink.org/RMSChannels}.
  \item Global Ionospheric Radio Observatory, “Ionogram Data Information,” \url{https://giro.uml.edu/ionogram-data.html}.
  \item Model Context Protocol, “Server Tools,” \url{https://modelcontextprotocol.io/specification/2025-06-18/server/tools}.
  \item Linkwise project, “HF Propagation Toolbox,” local project document: \path{reference/propagation-toolbox.md}.
  \item 8r4n, \emph{acq-latex-templates}, repository commit \texttt{fef160e}, \url{https://github.com/8r4n/acq-latex-templates}.
  \item Winlink, “FAQ, March 11, 2023,” channel-list maintenance guidance and Winlink Express update workflow, \url{https://winlink.org/content/winlink_faq_march_11_2023}.
  \item Winlink, “Winlink Express Quick Setup Guide,” locator and service-code configuration and channel-selector workflow, \url{https://winlink.org/sites/default/files/RMSE_FORMS/quick_setup_guide_for_winlink_sound_card_packet_for_vhf-uhf_on_windows_v1.3.pdf}.
  \item Winlink, “Glossary,” RMS/CMS and MID definitions, \url{https://winlink.org/glossary}.
  \item Winlink, “Overview of the Winlink Hybrid Network,” \url{https://winlink.org/content/overview_winlink_hybrid_network}.
  \item ITU-R, “Recommendation P.533: Method for the prediction of the performance of HF circuits,” \url{https://www.itu.int/rec/R-REC-P.533/en}.
  \item ITU-R, “Recommendation F.1487: Testing of HF modems with ionospheric channel simulators,” \url{https://www.itu.int/rec/R-REC-F.1487-0-200005-I/en}.
  \item GNU Radio, “Dynamic Channel Model block definition,” \url{https://github.com/gnuradio/gnuradio/blob/main/gr-channels/grc/channels_dynamic_channel_model.block.yml}.
  \item Jawatson, “VOACAPL Linux port,” \url{https://github.com/jawatson/voacapl}.
  \item Defence Science and Technology Group, “PHaRLAP HF ray-tracing toolbox” (optional future research), \url{https://www.dst.defence.gov.au/node/1870}.
  \item PlantUML, “Command line,” output formats and diagram generation, \url{https://plantuml.com/command-line}.
  \item KeepTrack, “KeepTrack API v4 documentation,” reentry endpoint, authentication, rate limits, and CC BY-NC 4.0 terms, \url{https://api.keeptrack.space/v4/docs}.
  \item Chiu, Yi-Chung, Loren C. Chang, and Jann-Yenq Liu, “The Ionospheric Impact of Artificial Space Object Reentry,” peer-reviewed abstract, Proceedings of the ION 2026 Pacific PNT Meeting, April 2026, \url{https://www.ion.org/publications/abstract.cfm?articleID=20596}.
  \item Linkwise project, “Software Design Description,” \texttt{docs/architecture/software-design.md}, proposed design baseline, \url{../architecture/software-design.md}.
  \item Wing, R., Gerding, M., Plane, J. M. C., et al., “Measurement of a lithium plume from the uncontrolled re-entry of a Falcon 9 rocket,” \emph{Communications Earth \& Environment} 7, 161 (2026), \url{https://doi.org/10.1038/s43247-025-03154-8}. Reports an upper-atmospheric lithium plume attributed to a specific re-entry; it does not measure HF/10 m propagation.
  \item Hamlib, “Ham radio control library for rigs, rotators, and amplifiers,” \url{https://github.com/Hamlib/Hamlib}.
  \item Winlink, “VARA Products Now Downloadable Here,” third-party software status, \url{https://winlink.org/content/vara_products_now_downloadable_here}.
  \item Model Context Protocol, “Rust SDK (rmcp),” \url{https://rust.sdk.modelcontextprotocol.io/}.
  \item RustAudio, “CPAL: Cross-Platform Audio Library,” \url{https://docs.rs/cpal/latest/cpal/}.
  \item Winlink, “Alternate Information,” client/modem configuration information including VARA interface ports, \url{https://winlink.org/content/alternate_information}.
  \item NOAA Space Weather Prediction Center, “Products,” machine-readable solar, geomagnetic, forecast, and warning products, \url{https://services.swpc.noaa.gov/products/}.
  \item NASA Community Coordinated Modeling Center, “DONKI,” event database, API, and research and prototyping limitations, \url{https://ccmc.gsfc.nasa.gov/tools/DONKI/}.
  \item NASA Earthdata CDDIS, “GNSS IGS Analysis Center Ionosphere Vertical Total Electron Content and GUARDIAN Near Real-Time TEC collections,” \url{https://cmr.earthdata.nasa.gov/search/site/collections/directory/CDDIS/gov.nasa.eosdis}.
  \item NOAA Space Weather Prediction Center, “SWPC Text Products,” including D-RAP and current geomagnetic forecast products, \url{https://services.swpc.noaa.gov/text/}.
  \item Budden, K. G., \emph{The Propagation of Radio Waves: The Theory of Radio Waves of Low Power in the Ionosphere and Magnetosphere}, Cambridge University Press, 1985, \url{https://doi.org/10.1017/CBO9780511564321}.
  \item Davies, K., \emph{Ionospheric Radio}, Peter Peregrinus / Institution of Electrical Engineers, 1990, \url{https://doi.org/10.1049/PBEW031E}.
  \item Hargreaves, J. K., \emph{The Solar-Terrestrial Environment: An Introduction to Geospace}, Cambridge University Press, 1992, \url{https://doi.org/10.1017/CBO9780511628924}.
  \item Hunsucker, R. D., and J. K. Hargreaves, \emph{The High-Latitude Ionosphere and its Effects on Radio Propagation}, Cambridge University Press, 2002, \url{https://doi.org/10.1017/CBO9780511535758}.
  \item Herv\'as, M., P. Bergad\`a, and R. M. Alsina-Pag\`es, “Ionospheric Narrowband and Wideband HF Soundings for Communications Purposes: A Review,” \emph{Sensors}, 20(9), 2486, 2020, \url{https://doi.org/10.3390/s20092486}.
  \item Ads, A. G., P. Bergad\`a, C. Vilella, J. R. Regu\'e, J. L. Pijoan, R. Bardaj\'i, and J. Mauricio, “A comprehensive sounding of the ionospheric HF radio link from Antarctica to Spain,” \emph{Radio Science}, 48(1), 1--12, 2012, \url{https://doi.org/10.1029/2012RS005074}.
  \item Walden, M. C., “Comparison of propagation predictions and measurements for midlatitude HF near-vertical incidence sky wave links at 5 MHz,” \emph{Radio Science}, 47, RS0L09, 2012, \url{https://doi.org/10.1029/2011RS004914}.
  \item Moskaleva, E. V., and N. Y. Zaalov, “A simplified HF radio channel forecasting model. Case study,” \emph{Advances in Space Research}, 69(6), 2477--2488, 2022, \url{https://doi.org/10.1016/j.asr.2021.12.047}.
  \item Zhu, Z., Q. Yu, Y. Shi, C. Yang, J. Fan, and J. Wang, “Exploring the Applicability of Three Popular Models for Predicting the Received Field Strength of HF Sky-Wave Propagation and Communication,” \emph{Space Weather}, 24, e2025SW004907, 2026, \url{https://doi.org/10.1029/2025SW004907}.
  \item Zawdie, K. A., D. P. Drob, D. E. Siskind, and C. Coker, “Calculating the absorption of HF radio waves in the ionosphere,” \emph{Radio Science}, 52, 767--783, 2017, \url{https://doi.org/10.1002/2017RS006256}.
  \item Frissell, N. A., et al., “Ionospheric Sounding Using Real-Time Amateur Radio Reporting Networks,” \emph{Space Weather}, 12, 651--656, 2014, \url{https://doi.org/10.1002/2014SW001132}.
  \item Frissell, N. A., et al., “First Observations of Large Scale Traveling Ionospheric Disturbances Using Automated Amateur Radio Receiving Networks,” \emph{Geophysical Research Letters}, 49, e2022GL097879, 2022, \url{https://doi.org/10.1029/2022GL097879}.
  \item Frissell, N. A., et al., “Heliophysics and amateur radio: citizen science collaborations for atmospheric, ionospheric, and space physics research and operations,” \emph{Frontiers in Astronomy and Space Sciences}, 10, 1184171, 2023, \url{https://doi.org/10.3389/fspas.2023.1184171}.
  \item Chartier, A. T., J. Steele, G. Sugar, D. R. Themens, S. K. Vines, and J. D. Huba, “Validating Ionospheric Models Against Technologically Relevant Metrics,” \emph{Space Weather}, 21, e2023SW003590, 2023, \url{https://doi.org/10.1029/2023SW003590}.
  \item Icom, “IC-7200,” product specifications and USB/CI-V capabilities, \url{https://www.icomjapan.com/lineup/products/IC-7200/}.
  \item Icom, “IC-7200 Instruction Manual,” official support/download page, \url{https://www.icomjapan.com/support/manual/1507/}.
  \item U.S. Department of Defense, “MIL-STD-188-141, Interoperability and Performance Standards for Medium and High Frequency Radio Systems,” ASSIST document record, \url{https://quicksearch.dla.mil/qsDocDetails.aspx?ident_number=67563}.
  \item Barrett Communications, “ALE 2G and 3G User Guide,” LQA exchange, averaging, decay, and threshold configuration, \url{https://www.barrettcommunications.com.au/Manuals/Barrett_ALE_User_Guide.pdf}.
  \item SoapySDR, “Device Class Reference,” device configuration and streaming API, \url{https://pothosware.github.io/SoapySDR/doxygen/latest/classSoapySDR_1_1Device.html}.
  \item GNU Radio, “Soapy,” source-block integration documentation, \url{https://wiki.gnuradio.org/index.php/Soapy}.
  \item\label{ref:modem-protocol} Saul St John (W9SSJ), \emph{hfmodem}, published VARA protocol notes, reference audio and measured waveform constants; revision \path{952f79245240132f4af209afbad1f4b2ba256613}, \url{https://github.com/sstjohn/hfmodem}. Protocol documentation and downloaded measurement sources retain their separate attribution/license metadata.
  \item Linkwise project, native modem implementation inventory and validation limits, \path{rust/docs/modem.md}; dated execution evidence in \path{lab-notes/lab-notes.tex}. The named Rust modules in Section~\ref{sec:modem-dsp} define the implementation facts described here.
\end{enumerate}

\end{document}
